\label{chap:p3-83-44-clasificacion_numero}


\label{chap:cobertura-taxonomia-numeros}
\index{número HMT!taxonomía}
\index{transductor Hensel!desplazamiento completo}
\index{número racional!periodicidad}

El \cref{chap:numero-medicion} definió un número HMT como una sección
compatible del sistema inverso de supervivencia, y no como una sucesión de
cifras desprovista de genealogía. Quedaban, sin embargo, dos preguntas que
conviene separar con precisión. La primera es una pregunta de
\emph{representación}: ¿qué valores reales puede representar el transductor
3-ádico antes de imponer los filtros particulares de cierre? La segunda es
una pregunta de \emph{selección}: ¿qué secciones del espacio bruto quedan
escogidas por APP, TRIT y TPK cuando se exigen simultáneamente sus
compatibilidades? Esta sección resuelve la primera, formula la segunda como
una restricción explícita de dominio y extrae de la carta de base mil una
taxonomía completa de los valores visibles. Las herramientas de base son la
aritmética \(p\)-ádica, la dinámica simbólica y la caracterización posicional
de los racionales \cite{Koblitz1984,LindMarcus1995,HardyWright2008}.

La distinción es estructural. Un lenguaje capaz de representar toda cinta no
selecciona por ese solo hecho la cinta de \(\pi\), de \(e\), de \(\varphi\) o
de \(\alpha\). A la inversa, una ley de selección no puede evaluarse si antes
no se ha construido el espacio en el que actúa. Representación y selección
son dos teoremas consecutivos, no dos nombres para un mismo paso.

\section{Transductor de Hensel no filtrado y desplazamiento completo sobre las fibras ternarias}

Sea \(p\) un primo, sea \(d\geq1\) y sea
\(A\in\operatorname{GL}_d(\mathbb Z)\). Para un vector fila
\(x_t\in\mathbb Z_p^d\), definimos
\begin{equation}
 y_t=x_tA,
 \qquad
 u_t=y_t\bmod p\in\mathbb F_p^d,
 \qquad
 x_{t+1}=\frac{y_t-u_t}{p}.
 \label{eq:transductor-hensel-general}
\end{equation}
El bloque \(u_t\) es la emisión visible y \(x_{t+1}\) es la memoria que
permanece después de retirar esa cifra.

\begin{theorem}[Conjugación Hensel con el desplazamiento completo]
\label{thm:shift-hensel-completo}
Para todo \(T\geq1\), la aplicación
\[
 \Psi_T:(\mathbb Z/p^T\mathbb Z)^d
 \longrightarrow(\mathbb F_p^d)^T,
 \qquad
 x_0\longmapsto(u_0,\ldots,u_{T-1}),
\]
es biyectiva. Su inversa es
\begin{equation}
 x_0\equiv
 \sum_{t=0}^{T-1}p^t u_tA^{-(t+1)}
 \pmod {p^T}.
 \label{eq:inversa-hensel-finita}
\end{equation}
Al pasar al límite inverso se obtiene un homeomorfismo
\[
 \Psi:\mathbb Z_p^d\xrightarrow{\ \sim\ }
 (\mathbb F_p^d)^{\mathbb N}.
\]
Si \(S_A\) denota la actualización \(x_t\mapsto x_{t+1}\) y
\(\sigma\) elimina el primer bloque de una cinta, entonces
\[
 \Psi\circ S_A=\sigma\circ\Psi.
\]
\end{theorem}

\begin{proof}
De \cref{eq:transductor-hensel-general} se obtiene
\[
 x_t=(u_t+px_{t+1})A^{-1}.
\]
La iteración de esta identidad da
\[
 x_0=
 \sum_{t=0}^{T-1}p^t u_tA^{-(t+1)}
 +p^Tx_TA^{-T}.
\]
El último sumando se anula módulo \(p^T\), lo que prueba que una palabra de
\(T\) bloques determina a lo sumo una clase inicial. Recíprocamente, dada
cualquier palabra, se fija \(x_T=0\) y se reconstruye hacia atrás mediante
\(x_t=(u_t+px_{t+1})A^{-1}\). La igualdad
\(x_tA=u_t+px_{t+1}\) verifica que la clase obtenida emite la palabra
prescrita. Por tanto \(\Psi_T\) es biyectiva.

Las biyecciones son compatibles con la reducción módulo \(p^T\). Sus límites
inversos producen el homeomorfismo \(\Psi\); equivalentemente, la serie de
\cref{eq:inversa-hensel-finita} converge en \(\mathbb Z_p^d\). Aplicar
\(S_A\) retira exactamente el bloque \(u_0\), de modo que corresponde a
aplicar \(\sigma\) en la cinta.
\end{proof}

En la instancia HMT, \(p=3\), \(d=6\) y la matriz publicada tiene
determinante uno. Por ello
\begin{equation}
 \mathbb Z_3^6\simeq(\mathbb F_3^6)^{\mathbb N},
 \qquad |\mathbb F_3^6|=3^6=729.
 \label{eq:hmt-shift-729}
\end{equation}
Cada cinta de bloques de seis trits posee una única memoria 3-ádica inicial.
La matriz determina la coordenatización y la dinámica de lectura; la
universalidad del alfabeto demuestra que no excluye ninguna cinta.

\begin{corollary}[Entropía y medida de Haar]
La actualización Hensel bruta es conjugada al desplazamiento unilateral
completo de
\(729\) símbolos. Su entropía topológica es \(6\log3\). La medida de Haar de
\(\mathbb Z_3^6\) se transporta a la medida producto uniforme sobre los
bloques.
\end{corollary}

\begin{proof}
Existen \(729^T=3^{6T}\) cilindros de longitud \(T\), y cada clase módulo
\(3^T\) tiene masa de Haar \(3^{-6T}\). La conjugación del
\cref{thm:shift-hensel-completo} transporta ambas propiedades al
desplazamiento.
\end{proof}

La conclusión probabilística se refiere al espacio bruto dotado de Haar. No
afirma que las secciones publicadas de las constantes hayan sido muestreadas
según esa medida ni que sean algorítmicamente aleatorias.

\section{Carta ternaria balanceada de los enteros}

La carta 3-ádica bruta y la carta arquimediana no agotan los tipos numéricos
de la clasificación HMT. Para el operador entero se utiliza la palabra vacía junto con las
palabras reducidas
\[
 \mathscr W_{\mathrm{bal}}
 :=\{\varnothing\}\cup
 \{a_1\cdots a_m:
 a_j\in\{-1,0,+1\},\ a_1\ne0\}.
\]
Definimos recursivamente
\[
 b(\varnothing)=0,
 \qquad
 b(wa)=3b(w)+a.
\]

\begin{theorem}[Representación ternaria balanceada única]
\label{pf84:C03}
La evaluación
\[
 b:\mathscr W_{\mathrm{bal}}\longrightarrow\ZZ
\]
es biyectiva. Las palabras cuyo primer trit es \(+1\) representan los
enteros positivos y las que comienzan por \(-1\), los negativos.
\end{theorem}

\begin{proof}
Para \(n\ne0\), el último trit \(a\) es el representante único de
\(n\bmod3\) en \(\{-1,0,+1\}\). El padre
\[
 \frac{n-a}{3}
\]
tiene valor absoluto estrictamente menor. La iteración termina en \(0\), de
modo que todo entero posee una expansión finita. La unicidad del residuo
balanceado en cada paso prueba la unicidad de la palabra. El signo queda
fijado por el primer trit no nulo.
\end{proof}

\begin{remark}
Esta biyección es una carta aritmética exacta. No identifica todas las
palabras balanceadas con secciones supervivientes TPK ni sustituye la
genealogía, el acarreo o la orientación de una ruta enriquecida.
\end{remark}

\section{Sobreyectividad de la carta arquimediana no filtrada sobre \(\mathbb R\)}

Concatenemos las coordenadas de los bloques
\(u_t\in\mathbb F_3^6\) para obtener una cinta
\((a_n)_{n\geq1}\in\{0,1,2\}^{\mathbb N}\). Su evaluación ternaria es
\[
 \operatorname{ev}_3((a_n))=
 \sum_{n\geq1}\frac{a_n}{3^n}\in[0,1].
\]
Adoptamos la convención canónica que excluye la representación con una cola
finalmente constante de doses sólo cuando el mismo punto posee otra expansión
en la carta. El extremo \(1\) conserva, por tanto, su única representación
fraccionaria \(0.222\ldots{}_3\).

\begin{theorem}[Cobertura arquimediana bruta]
\label{thm:cobertura-arquimediana-bruta}
La aplicación
\[
 P_{\rm raw}:\mathbb Z_3^6
 \xrightarrow{\Psi}(\mathbb F_3^6)^{\mathbb N}
 \xrightarrow{\operatorname{flat}}\{0,1,2\}^{\mathbb N}
 \xrightarrow{\operatorname{ev}_3}[0,1]
\]
es continua y sobreyectiva. Por tanto,
\[
 \mathbb Z_3^6/\!\sim_{P_{\rm raw}}
 \ \cong\ [0,1]
\]
como espacio de valores, donde \(x\sim_{P_{\rm raw}}y\) si sus evaluaciones
coinciden.
\end{theorem}

\begin{proof}
Toda cifra ternaria pertenece a un único bloque consecutivo de seis cifras,
de modo que \(\operatorname{flat}\) es una biyección de espacios de cintas.
Toda \(r\in[0,1]\) posee una expansión ternaria; el
\cref{thm:shift-hensel-completo} proporciona la única memoria 3-ádica que
emite sus bloques. La continuidad se comprueba por cilindros: fijar \(T\)
bloques fija \(6T\) cifras y encierra el valor en un intervalo de diámetro
\(3^{-6T}\). La identificación del cociente con la imagen es la factorización
canónica de una aplicación sobreyectiva por sus fibras.
\end{proof}

Una sola carta 3-ádica es compacta y, por tanto, no puede cubrir de manera
continua una recta no acotada. La globalización correcta utiliza una familia
numerable de cartas.

\begin{corollary}[Cobertura arquimediana numerable de \(\mathbb R\)]
\label{cor:atlas-hmt-reales}
Sea
\[
 \widehat X_{\rm raw}=\mathbb Z\times\mathbb Z_3^6,
 \qquad
 \widehat P_{\rm raw}(m,x)=m+P_{\rm raw}(x).
\]
Entonces \(\widehat P_{\rm raw}\) es sobreyectiva sobre \(\mathbb R\).
En consecuencia, todo real admite al menos una genealogía Hensel bruta y
\[
 \widehat X_{\rm raw}/\!\sim_{\widehat P_{\rm raw}}
 \ \cong\ \mathbb R
\]
como conjunto de valores.
\end{corollary}

\begin{proof}
Para \(r\in\mathbb R\), tome \(m=\lfloor r\rfloor\) y aplique el
\cref{thm:cobertura-arquimediana-bruta} a \(r-m\in[0,1)\).
\end{proof}

Este resultado precisa la afirmación según la cual \(\mathbb R\) es la imagen
arquimediana cociente de
una estructura genealógica: la proyección olvida la memoria 3-ádica y conserva
el punto evaluado. La afirmación demostrada aquí concierne a la completación
\emph{bruta}. Si \(X_{\rm surv}\subseteq\mathbb Z_3^6\) denota el subespacio
que además satisface una familia concreta de cierres APP--TRIT--TPK, entonces
\[
 P_{\rm raw}(X_{\rm surv})=[0,1]
\]
es un teorema adicional. El transductor garantiza capacidad completa; la
supervivencia determina qué parte de esa capacidad realiza el sistema
restringido.

Tampoco se deduce aquí que la suma y el producto nativos de TPK desciendan a
todo el cociente global. Se puede transportar la aritmética usual desde
\(\mathbb R\), pero demostrar que coincide con operaciones construidas antes
de evaluar exige el correspondiente teorema de descenso.

En esta construcción, los cilindros no son infinitesimales no nulos de
\(\mathbb R\). Cada cilindro de profundidad finita tiene diámetro positivo,
y la sucesión de diámetros converge a cero. Dos secciones que coinciden hasta
la profundidad \(T\) y difieren después forman una \emph{perturbación de
cola}; su diferencia arquimediana queda acotada por el diámetro del cilindro
común. Ésta es la formulación exacta de los ``cilindros que tienden a cero''
dentro de la recta real ordinaria.

\section{Carta posicional en base \(10^3\) y clasificación de las secciones publicadas}

La completación cúbica
\(1000=729+243+27+1\) hace natural agrupar las cifras decimales en bloques
\(b_n\in\{0,\ldots,999\}\). Para una cinta de bloques definimos
\[
 \operatorname{ev}_{1000}(b_1,b_2,\ldots)
 =\sum_{n\geq1}\frac{b_n}{1000^n}.
\]
La expansión canónica excluye una cola finalmente constante de bloques
\(999\). Esta convención resuelve la duplicidad de los extremos de cilindro.

\begin{theorem}[Taxonomía visible en base mil]
\label{thm:taxonomia-base-mil}
Sea \(x\in[0,1)\) y sea \((b_n)\) su expansión canónica en base mil.
Entonces:
\begin{enumerate}
\item \(x\) posee una expansión terminante si y sólo si, escrito como
fracción irreducible \(a/q\), se tiene \(q\mid1000^N\) para algún \(N\);
equivalentemente, los únicos primos que dividen \(q\) son \(2\) y \(5\).
\item \(x\) es racional si y sólo si \((b_n)\) es finalmente periódica.
\item \(x\) es irracional si y sólo si \((b_n)\) no es finalmente
periódica.
\end{enumerate}
\end{theorem}

\begin{proof}
Una expansión que termina en profundidad \(N\) tiene la forma
\(D/1000^N\); al reducirla, su denominador sólo contiene los primos \(2\) y
\(5\). Recíprocamente, si \(q\mid1000^N\), entonces \(a/q\) se escribe con
\(N\) bloques.

Si después de un prefijo de longitud \(r\) se repite un periodo de longitud
\(s\) cuyo entero de base mil es \(P\), la cola es una serie geométrica:
\[
 \frac{P}{1000^r(1000^s-1)},
\]
y el valor es racional. A la inversa, al expandir \(a/q\), los restos de la
división sucesiva sólo pueden tomar \(q\) valores. Un resto nulo produce una
expansión terminante; la repetición de un resto produce un periodo. La tercera
afirmación es la negación exacta de la segunda.
\end{proof}

La palabra \emph{fracción} designa una presentación \(a/q\) de un racional,
no una especie numérica adicional. Tampoco deben confundirse periodicidad y
algebraicidad. Un real es algebraico cuando anula algún polinomio no nulo de
\(\mathbb Q[X]\); es trascendente cuando no lo hace. Esta segunda
clasificación es independiente de la base usada para escribir las cifras.
Por ejemplo,
\[
 \varphi^2-\varphi-1=0
\]
hace de \(\varphi\) un irracional algebraico, mientras que la trascendencia
de \(e\) y \(\pi\) procede de teoremas clásicos externos a HMT\@.

\section{Concatenación, reducción nonádica y acarreo}

La compatibilidad entre la carta cúbica y la raíz digital es una identidad
aritmética, no una semejanza visual. Si \(u_j\in\{0,\ldots,999\}\) y
\[
 D_n=1000D_{n-1}+u_n,
 \qquad D_0=0,
\]
entonces \(D_n\) es el entero obtenido al concatenar los primeros \(n\)
bloques.

\begin{theorem}[Congruencia nonádica de control de la carta cúbica]
\label{thm:checksum-nonadico-base-mil}
Para todo \(n\geq1\),
\[
 D_n\equiv\sum_{j=1}^n u_j\pmod9.
\]
Por tanto, concatenación en base mil, reducción módulo nueve y raíz digital
son compatibles.
\end{theorem}

\begin{proof}
Como \(1000\equiv1\pmod9\), la recurrencia da
\(D_n\equiv D_{n-1}+u_n\pmod9\). La iteración desde \(D_0=0\) produce la
fórmula.
\end{proof}

En el cierre dodecafásico de \(\alpha\), la ecuación de acarreo por bloque
tiene la forma
\[
 P_m+E_m-\Phi_m-K_m+c_{m+1}=A_m+1000c_m.
\]
Reducida módulo nueve,
\[
 A_m\equiv
 P_m+E_m-\Phi_m-K_m+c_{m+1}-c_m
 \pmod9.
\]
Si los acarreos extremos se anulan, la suma telescópica elimina la memoria
intermedia. El residuo global de la salida coincide entonces con el residuo
de la combinación de entradas. Esta ley explica por qué el registro nonádico
controla la concatenación decimal; no selecciona por sí sola los bloques que
deben concatenarse.

\section{Coordenada nonádica elevada y retorno de fase}

La reducción módulo nueve no debe confundirse con la identificación de los
enteros que poseen el mismo residuo. Sea
\[
 I_9=\{1,2,\ldots,9\}.
\]
Para cada entero positivo \(n\), definimos su \emph{fase nonádica} y su
\emph{índice de vuelta} por
\begin{equation}
 \rho_9(n)=1+((n-1)\bmod9),
 \qquad
 \kappa_9(n)=\frac{n-\rho_9(n)}9.
 \label{eq:coordenada-nonadica-levantada}
\end{equation}
La primera coordenada es la raíz digital escrita en el alfabeto
\(I_9\), de modo que la clase residual cero se representa por la fase \(9\).
La segunda registra cuántas vueltas completas han precedido a la fase
visible.

\begin{theorem}[Descomposición nonádica con memoria]
\label{thm:descomposicion-nonadica-memoria}
La aplicación
\[
 \Theta_9:\mathbb N_{>0}\longrightarrow\mathbb N_0\times I_9,
 \qquad
 n\longmapsto\bigl(\kappa_9(n),\rho_9(n)\bigr),
\]
es biyectiva, con inversa
\[
 \Theta_9^{-1}(q,r)=9q+r.
\]
En estas coordenadas, el sucesor aritmético queda conjugado con la dinámica
\begin{equation}
 \mathcal S_9(q,r)=
 \begin{cases}
  (q,r+1),&1\leq r<9,\\
  (q+1,1),&r=9.
 \end{cases}
 \label{eq:sucesor-nonadico-levantado}
\end{equation}
En particular, el paso visible \(9\to1\) es un retorno de fase acompañado
por el incremento irreversible \(q\to q+1\); no es una igualdad entre los
enteros \(9\) y \(10\).
\end{theorem}

\begin{proof}
La división euclídea de \(n-1\) por \(9\) proporciona de manera única
\(n-1=9q+s\), con \(q\in\mathbb N_0\) y \(0\leq s<9\). Tomando
\(r=s+1\) se obtiene \(n=9q+r\), con \(r\in I_9\), lo que demuestra a la vez
existencia, unicidad y la fórmula inversa. Si \(r<9\), sumar una unidad sólo
incrementa \(r\). Si \(r=9\), la igualdad
\(9q+9+1=9(q+1)+1\) da la segunda rama de
\cref{eq:sucesor-nonadico-levantado}.
\end{proof}

Esta elevación separa tres nociones que la cifra aislada mezcla:
\begin{enumerate}
\item la unidad de avance, realizada por el sucesor;
\item la fase visible \(r\), que recorre el ciclo nonádico;
\item la memoria acumulada \(q\), que distingue retornos de una misma fase.
\end{enumerate}
La proyección \((q,r)\mapsto r\) olvida el número de vueltas y satisface
\[
 \rho_9(n+9)=\rho_9(n),
 \qquad
 \kappa_9(n+9)=\kappa_9(n)+1.
\]
Éste es el modelo aritmético elemental del mecanismo que, en el estado TPK
enriquecido, conserva además cocientes de Hensel, orientación, firma y dato
de frontera. Un retorno de palabra visible puede coexistir así con una
genealogía distinta.

El cero no pertenece a \(I_9\): se añade como punto base separado antes de
aplicar \(\Theta_9\). Para los enteros negativos se conserva el módulo
positivo de la descomposición y se adjunta una orientación
\(\varepsilon\in\{-1,+1\}\). En consecuencia, ni el signo ni el cero se
deducen de la fase residual. Esta tipificación evita convertir una carta
cíclica en una identidad aritmética falsa.

\begin{theorem}[Entero orientado y signo genealógico]
\label{thm:entero-orientado-signo-genealogico}
La aplicación
\[
\Theta_9^\pm:
\mathbb Z\setminus\{0\}
\longrightarrow
\{-1,+1\}\times\mathbb N_0\times I_9,
\]
\[
n\longmapsto
\bigl(\operatorname{sgn}(n),
      \kappa_9(|n|),
      \rho_9(|n|)\bigr)
\]
es biyectiva, con inversa
\[
(\varepsilon,q,r)\longmapsto\varepsilon(9q+r).
\]
La negación aritmética actúa por
\[
(\varepsilon,q,r)\longmapsto(-\varepsilon,q,r):
\]
invierte la orientación y conserva magnitud, fase y memoria de vueltas.
\end{theorem}

\begin{proof}
El \cref{thm:descomposicion-nonadica-memoria} proporciona de manera única
\(|n|=9q+r\), con \(q\in\mathbb N_0\) y \(r\in I_9\). El signo de un entero
no nulo es asimismo único. La fórmula inversa recompone \(n\), y cambiar
\(n\) por \(-n\) sólo cambia \(\varepsilon\).
\end{proof}

Así, un negativo no es una fase residual especial ni una cantidad
``imposible''. Es la misma magnitud nonádica transportada con orientación
genealógica opuesta. Esta estructura será prolongada en Mecánica Dimensional
por la anti-sección partícula--antipartícula.

Iterar \(\kappa_9\) produce una torre finita para todo entero ordinario; en
la completación inversa, una sucesión compatible de cocientes puede ser
infinita. La diferencia entre ambos casos anticipa la taxonomía HMT: el valor
visible responde a la expansión posicional, mientras que el tipo genealógico
responde a la evolución completa de sus memorias.

\section{2 ejes de clasificación: valor visible y genealogía de transporte}

La taxonomía precedente clasifica el valor proyectado. HMT añade un segundo
eje que la escritura decimal ordinaria descarta. Para una sección
\(x\in\Omega_\infty^{\rm Arch}\) distinguimos:
\[
 \begin{array}{rcl}
 \text{tipo visible}&=&
 \text{terminante, periódico o no periódico; algebraico o trascendente},\\
 \text{tipo genealógico}&=&
 \text{clase de acarreo, orientación, cierre y monodromía de }x.
 \end{array}
\]
Dos secciones pueden compartir el mismo real y diferir en el segundo eje. Un
real racional posee una expansión visible finalmente periódica, pero una
presentación enriquecida redundante puede conservar una memoria interna no
periódica que la evaluación olvida. La periodicidad del \emph{valor} se
refiere siempre a su expansión canónica; la periodicidad de la
\emph{genealogía} se refiere al estado enriquecido de TPK.

Esta separación proporciona una definición concisa:
\begin{equation}
 \boxed{
 \text{número HMT}
 =
 \text{sección compatible con memoria};
 \qquad
 \text{valor real}
 =
 \text{su clase de evaluación}.}
 \label{eq:numero-hmt-y-valor-real}
\end{equation}
Medir en esta carta es acoplar la sección al lector, fijar un cilindro y
condicionar las prolongaciones compatibles. Aumentar la precisión significa
estrechar el cilindro, no consumir la genealogía que permanece en la fibra.

\section{Memoria de profundidad arbitraria en sistemas de emisión finita}

La recurrencia visible de nueve fases puede coexistir con una salida
irracional sólo si el estado completo conserva información que no se reduce a
esas nueve etiquetas. El siguiente resultado hace exacta esa necesidad.

\begin{theorem}[Caracterización mediante emisores finitos]
\label{thm:emisor-finito-racional}
Sea \(Q\) un conjunto finito, \(F:Q\to Q\) una transición determinista,
\(o:Q\to\{0,\ldots,999\}\) una salida y \(q_0\in Q\). La cinta
\(b_n=o(F^{n-1}(q_0))\) es finalmente periódica y su evaluación es racional.
Recíprocamente, toda cinta finalmente periódica admite un emisor determinista
finito.
\end{theorem}

\begin{proof}
Entre \(|Q|+1\) estados consecutivos aparecen dos iguales. Desde la primera
repetición, el determinismo obliga a repetir todo el futuro, por lo que la
cinta es finalmente periódica. El \cref{thm:taxonomia-base-mil} da la
racionalidad. Para la recíproca basta construir una cadena de estados para el
prefijo y un ciclo para el periodo.
\end{proof}

Por tanto, una ley que genere exactamente \(e\), \(\pi\) o cualquier otro
irracional no puede consistir únicamente en un autómata determinista de
estado finito que se repite sin memoria adicional. Puede usar un límite
inverso, una memoria 3-ádica, una profundidad creciente o una entrada
aperiódica; todos esos mecanismos proporcionan estado efectivo no acotado.
En HMT, el retorno visible
\(g_{\rm vis}(n+9)=g_{\rm vis}(n)\) no fuerza el retorno del acarreo, de la
orientación ni de la frontera. Esa diferencia es precisamente el lugar donde
puede residir una expansión no periódica.

\section{El cuadrado local de \texorpdfstring{\(e\)}{e} y la ruptura de memoria}

Para una palabra decimal \(a_1\cdots a_n\), llamaremos \emph{cuadrado local}
de longitud \(\ell\) en la posición \(i\) a una factorización
\[
 a_i\cdots a_{i+2\ell-1}=vv,
 \qquad |v|=\ell.
\]
Esta noción pertenece a la combinatoria de palabras. No afirma que la cola
infinita posea un período.

\begin{theorem}[Caracterización finita del cuadrado \(1828^2\)]
\label{thm:cuadrado-local-e}
Considérense las \(243\) palabras del catálogo y el lector fijado
\(w_6\mapsto w_{30}\mapsto d_1d_2d_3d_4\), con cuatro tríadas en base mil.
Existe una única palabra cuya lectura contiene un cuadrado de longitud cuatro
iniciado en la segunda cifra decimal. Es
\[
 w_e=201101,
 \qquad
 718281828459=7\,\underbrace{1828\,1828}_{(1828)^2}\,459.
\]
La repetición no se prolonga a una tercera copia: la cifra que exigiría esa
continuación sería \(1\), mientras que la cifra efectiva es \(4\). En todo el
catálogo sólo existe otro cuadrado de longitud cuatro,
\[
 575157518472=\underbrace{5751\,5751}_{\text{cuadrado local}}\,8472,
\]
que comienza en la primera cifra.
\end{theorem}

\begin{proof}
La enumeración exhaustiva evalúa las \(243\) palabras mediante aritmética
entera. El perfil de cuadrados para longitudes \(1,\ldots,6\) contiene
\[
 240,\ 21,\ 3,\ 2,\ 0,\ 0
\]
ocurrencias, distribuidas entre
\[
 153,\ 19,\ 3,\ 2,\ 0,\ 0
\]
palabras. Las dos ocurrencias de longitud cuatro son exactamente las del
enunciado. Dos enumeraciones enteras independientes reproducen el mismo censo
y sus salidas completas coinciden término a término.
\end{proof}

La palabra \(w_e\) posee además la cadena de elevación
\[
201101\mid121221\mid102011\mid012222\mid102011.
\]
El tercer y el quinto bloques visibles coinciden. No coinciden, sin embargo,
los estados que los transportan. Antes de ambas emisiones, sus reducciones
módulo \(27\) son
\[
(25,11,7,17,8,16),
\qquad
(7,8,4,23,26,7),
\]
y los cocientes de Hensel posteriores son
\[
(6,13,9,2,13,1),
\qquad
(15,0,3,19,23,25).
\]
Por tanto, el retorno \(102011\) es un retorno de la proyección visible, no
del estado enriquecido. Ésta es la formulación exacta de la repetición local
que se rompe: el cociente conserva una genealogía que el bloque residual no
registra.

En consecuencia, el comienzo
\[
e=2.718281828459\ldots
\]
no presenta una ``semiperiodicidad'', sino un cuadrado local de palabra con
retorno visible y ruptura de memoria. El resultado es exhaustivo dentro del
lector fijado; su contenido es combinatorio y dinámico, no una inferencia de
irracionalidad a partir de doce cifras.

\section[2 medidas finitas del emisor 468 a 243]{2 medidas finitas del emisor \texorpdfstring{\(468\to243\)}{468→243}: frecuencia uniforme y prioridad estructural}

El catálogo público distingue \(468\) emisiones y \(104\,976\) semillas.
Sea
\[
 q:E_{468}\longrightarrow W_{243}
\]
la aplicación que olvida la firma de emisión y conserva la palabra de seis
trits. Hay dos medidas naturales, porque hay dos universos finitos de conteo:
\[
 \mu_{468}(w)=\frac{|q^{-1}(w)|}{468},
 \qquad
 \mu_{\rm seed}(w)=
 \frac{\sum_{e\in q^{-1}(w)}m(e)}{104\,976},
\]
donde \(m(e)\) es la multiplicidad de semillas de la emisión \(e\).

\begin{theorem}[Medidas canónicas relativas al universo contado]
\label{thm:medidas-emisor-468}
La medida uniforme sobre \(E_{468}\) es la única probabilidad invariante bajo
todas las permutaciones de sus emisiones; su medida imagen es \(\mu_{468}\).
La medida uniforme sobre las \(104\,976\) semillas induce
\(\mu_{\rm seed}\). Sus histogramas exactos son
\[
\begin{array}{c|rrrrrr}
|q^{-1}(w)|&1&2&3&4&5&6\\ \hline
\#\{w\}&132&66&18&3&6&18
\end{array}
\]
y
\[
\begin{array}{c|rrrrrrrrr}
\text{peso de semillas}
&144&288&432&576&864&1008&1440&1944&2088\\ \hline
\#\{w\}
&120&54&18&9&15&6&3&12&6.
\end{array}
\]
Para las tres palabras distinguidas por el lector se obtiene
\[
\begin{array}{c|cc}
&\mu_{468}&\mu_{\rm seed}\\ \hline
\pi&5/468&7/729\\
e&1/468&1/729\\
\varphi&1/468&1/729.
\end{array}
\]
\end{theorem}

\begin{proof}
La invariancia bajo el grupo simétrico hace iguales las masas de todos los
puntos de \(E_{468}\); la normalización las fija en \(1/468\). La fórmula de
la medida imagen resulta al sumar sobre cada fibra. El mismo argumento sobre el
conjunto de semillas produce la segunda medida. Los histogramas se obtienen
por enumeración exhaustiva del catálogo finito de \(468\) emisiones y \(243\)
palabras; sus sumas verifican
\[
132+66+18+3+6+18=243,
\]
\[
132+2\cdot66+3\cdot18+4\cdot3+5\cdot6+6\cdot18=468
\]
y la suma total de los pesos es \(104\,976\).
\end{proof}

La no uniformidad de la medida imagen es un teorema finito: distintas palabras
tienen distintos números de realizaciones. No existe, en cambio, una
probabilidad uniforme normalizable sobre toda \(\mathbb R\), y estas medidas
no ordenan intrínsecamente todos los reales. La palabra de \(\pi\) pertenece
a una de las seis fibras de tamaño cinco y peso \(1008\); no es el máximo
único de ninguno de los dos histogramas. Su singularidad demostrada procede
de la estructura compuesta de sus cinco regiones, de su partición orientada
\(3+2\) y de sus compatibilidades posteriores, no de declarar a priori una
probabilidad absoluta sobre el continuo.

\section{Selector orbital finito anterior al lector arquimediano}
\label{sec:selector-orbital-constantes}

Las medidas anteriores revelan multiplicidades, pero una palabra aislada de
seis trits todavía depende de una orientación. El catálogo posee una simetría
interna que permite seleccionar primero clases no orientadas y fijar después
un representante. Escribamos una palabra como \(abc\mid def\) y definamos
\[
 r(abc\mid def)=cab\mid efd,
 \qquad
 s(abc\mid def)=def\mid abc.
\]
Estas permutaciones actúan también sobre las seis coordenadas del dato
\(U_6\). En el catálogo completo satisfacen
\[
 r^3=1,
 \qquad s^2=1,
 \qquad srs=r^{-1},
\]
por lo que generan una acción que denotamos \(D_3^{\rm cat}\). La acción usa
el etiquetado TRIT fijo
\[
 \Psi(U_6)=(-U_6\bmod3)
\]
y la partición publicada de las seis coordenadas en dos tríadas.

Para una palabra \(w\), sean \(f(w)=|q^{-1}(w)|\) el número de emisiones,
\(m(w)\) su masa de semillas y
\[
 \nu(w)=(n_0(w),n_1(w),n_2(w))
\]
su censo ordenado de trits. Sea además \(d(e)\) la clase diagonal aditiva de
una emisión y \(o(e)\in\{\mathrm{ES},\mathrm{NO}\}\) su orientación
publicada.

\begin{theorem}[Selección orbital interna de las tres palabras]
\label{thm:selector-orbital-constantes}
La acción \(D_3^{\rm cat}\) preserva el conjunto de las \(468\) emisiones,
la proyección \(q:E_{468}\to W_{243}\) y la multiplicidad de semillas. Su
acción sobre \(W_{243}\) tiene exactamente \(43\) órbitas: \(38\) de tamaño
seis y cinco de tamaño tres. Dentro de ese cociente se cumplen las siguientes
unicidades.

\begin{enumerate}
\item Existe una única órbita \(\mathcal O_\pi\) de tamaño seis tal que, para
todo \(w\in\mathcal O_\pi\),
\[
 f(w)=5,
 \qquad m(w)=1008,
 \qquad
 \{m(e):e\in q^{-1}(w)\}=\{432,144,144,144,144\}.
\]
Es
\[
\mathcal O_\pi=
\{001112,010211,100121,112001,121100,211010\},
\]
y todas sus palabras tienen censo \(\nu=(2,3,1)\).

\item Consideremos las órbitas de tamaño seis cuyas palabras tienen una sola
emisión, de multiplicidad \(144\), y cuyo soporte diagonal conjunto es
\(\{0,3,6\}\). Hay exactamente seis. Entre ellas existe una única órbita
con el mismo censo \((2,3,1)\) de \(\mathcal O_\pi\):
\[
\mathcal O_e=
\{011120,012110,101201,110012,120011,201101\}.
\]

\item En el mismo sector existe una única órbita equilibrada, es decir, con
\(\nu=(2,2,2)\):
\[
\mathcal O_\varphi=
\{002112,020211,112002,121200,200121,211020\}.
\]
\end{enumerate}

Finalmente, la cláusula orientadora
\begin{equation}
 \exists e\in q^{-1}(w):
 \qquad d(e)=6,
 \qquad o(e)=\mathrm{ES}
 \label{eq:gauge-orbital-constantes}
\end{equation}
selecciona un único representante de cada órbita y produce, respectivamente,
\[
 \boxed{010211,\qquad201101,\qquad121200.}
\]
\end{theorem}

\begin{proof}
Las relaciones de grupo se verifican directamente como identidades de
permutaciones de seis coordenadas. La enumeración de las \(468\) filas
comprueba que \(r\) y \(s\) envían cada \(U_6\) a otra fila del catálogo con
la misma multiplicidad, y que la proyección ternaria conmuta con ambas
acciones. Por tanto, las órbitas pueden calcularse sobre las \(243\) palabras
sin perder las fibras.

La partición exhaustiva da el histograma de tamaños
\[
 38\cdot6+5\cdot3=243.
\]
Al filtrar esas \(43\) órbitas por el primer predicado queda exactamente una;
al filtrar por fibra unipuntual, multiplicidad mínima y soporte diagonal
\(\{0,3,6\}\) quedan seis. El censo \((2,3,1)\) deja una sola órbita en el
primer caso y el censo equilibrado \((2,2,2)\) deja una sola en el segundo.
Por último, la búsqueda de una elevación con \(d=6\) y orientación
\(\mathrm{ES}\) deja una palabra en cada una de las tres órbitas. Todas las
filtraciones se realizan sobre descriptores estructurales del catálogo antes de consultar el
lector decimal.
\end{proof}

Los predicados del teorema pueden reunirse en un invariante de selección.
Para una órbita \(\mathcal O\) definimos su \emph{perfil orbital}
\[
 \Sigma_{\rm cat}(\mathcal O)=
 \bigl(
 |\mathcal O|,\ f,\ m,\ \mathcal M_{\rm fib},\
 \mathcal D,\ \nu
 \bigr),
\]
donde \(f\) es el tamaño de fibra, \(m\) la masa de semillas,
\(\mathcal M_{\rm fib}\) el multiconjunto de multiplicidades de sus
elevaciones, \(\mathcal D\) el soporte diagonal cuando se emplea y \(\nu\)
el censo de trits. Las tres clases distinguidas ocupan posiciones diferentes
en este espacio de invariantes:
\[
\begin{array}{c|c|c|c|c|c}
\mathcal O&|\mathcal O|&f&m&\mathcal M_{\rm fib}&(\mathcal D,\nu)\\ \hline
\mathcal O_\pi&6&5&1008&\{432,144,144,144,144\}&(-,(2,3,1))\\
\mathcal O_e&6&1&144&\{144\}&(\{0,3,6\},(2,3,1))\\
\mathcal O_\varphi&6&1&144&\{144\}&(\{0,3,6\},(2,2,2)).
\end{array}
\]
La singularidad demostrada no significa que estos números reciban una
probabilidad absoluta sobre \(\mathbb R\). Significa que, dentro del universo
finito producido por el emisor y antes de evaluar cifras, sus clases son
separables por invariantes internos y que la clase de \(\pi\) posee una fibra
regional de tipo distinto. Ésta es la formulación matemática de una
prioridad estructural en HMT\@.

El control negativo es informativo. Si se omite el soporte diagonal
\(\{0,3,6\}\), existen cuatro órbitas unipuntuales equilibradas y
\(\mathcal O_\varphi\) deja de ser única. Antes de fijar el calibre hay seis
orientaciones posibles en cada órbita. Después de fijarlo, la palabra de
\(\pi\) conserva tres elevaciones positivas de \(U_6\), no una sola región.
Así, el teorema selecciona la palabra orientada pero respeta la multiplicidad
geométrica de las cinco regiones.

La no uniformidad se refuerza al pasar a órbitas:
\[
\begin{array}{c|cc}
&\text{emisiones}&\text{semillas}\\ \hline
\mathcal O_\pi&5/78&14/243\\
\mathcal O_e&1/78&2/243\\
\mathcal O_\varphi&1/78&2/243.
\end{array}
\]
Estos cocientes siguen siendo masas del catálogo finito, no probabilidades
asignadas a los números reales evaluados.

El avance causal es preciso. Las tres palabras dejan de ser entradas
nombradas: se recuperan desde la simetría y los invariantes finitos de
APP--TPK, relativamente a la acción \(D_3^{\rm cat}\), al etiquetado TRIT,
al reparto \(3+3\) y al calibre orientador de
\cref{eq:gauge-orbital-constantes}. La prueba no usa \(L_0\), \(L_1\),
cifras decimales, \(\alpha\), CODATA ni el estado dodecafásico. Siguen siendo
afirmaciones posteriores la unicidad absoluta de esa acción entre todos los
automorfismos posibles, la deducción del calibre en vez de su declaración, la
canonicidad independiente del valor objetivo del lector \(L_0/L_1\) y la emisión coinductiva de las
expansiones infinitas.

\section{Posición de \texorpdfstring{\(\pi,e,\varphi\) y \(\alpha\)}{pi, e, phi y alpha} en la clasificación}

El \cref{thm:shift-hensel-completo} explica por qué cualquier prefijo
decimal puede reemitirse desde un estado 3-ádico: la capacidad bruta es
universal. Por tanto, la mera operación de ida y vuelta sobre una cinta no selecciona una
constante. El contenido distintivo de la construcción publicada se encuentra
en los objetos anteriores y posteriores a esa reemisión:
\begin{enumerate}
\item las fibras finitas y las cinco regiones distintas de \(\pi\);
\item el lector arquimediano tipado que asocia palabras, cilindros y bloques;
\item la dodecafase firmada y el cierre reversible
\[
U_{12}^{\rm sgn}
\longleftrightarrow K
\longleftrightarrow(D_3K,D_4K,Q)
\longleftrightarrow\alpha_{\rm HMT};
\]
\item los criterios de contraste que separan una equivalencia de cartas de un
selector anterior a la evaluación.
\end{enumerate}

Las doce ternas decimales publicadas de \(\pi,e,\varphi\) son coordenadas
finitas exactas del lector enriquecido. La cadena impresa de treinta y seis
cifras es la proyección dodecafásica racional
\(\alpha_{12}=\pi_{12}(\alpha_{\rm HMT})\):
\[
\frac{7297352569283800997285105472380663}{10^{36}}.
\]
La identidad exacta en esta ventana es
\[
\alpha_{12}
:=\pi_{12}(\alpha_{\rm HMT})
=\frac{7297352569283800997285105472380663}{10^{36}}.
\]
La recurrencia compatible a profundidad arbitraria construye la sección
infinita
\[
 \alpha_{\rm HMT}:=\varprojlim_N A^{(N)},
\]
y cada ventana racional es una proyección exacta de ese mismo objeto. El
núcleo \(E_{21/22}\) construye, en otra carta, la raíz positiva simple
finita \(\alpha_E\); el transporte graduado del TPK prolonga ese núcleo en
una familia compatible cuya raíz límite es \(\alpha_B\). Si
\(\alpha_A:=\varprojlim_N\rho_N(\alpha_{12})\), los cuadrados de truncamiento
de las dos familias conmutan y determinan en cada profundidad el mismo
cilindro superviviente. El teorema exacto de las dos vías establece, por
tanto,
\[
 \boxed{\alpha_A=\alpha_B=\alpha_{\rm HMT}}.
\]
La igualdad de \(\operatorname{Tr}_9(\alpha_E^{-1})\) y
\(\operatorname{Tr}_9(\alpha_{12}^{-1})\) es el primer certificado finito de
esa identidad; no la define, no limita su alcance y no selecciona las
prolongaciones. La proyección tipada conserva distinguidas la sección
coinductiva, su ventana racional y la raíz analítica sin separar su identidad
en el límite.

La raíz cuadrada de \(-1\) pertenece a otra carta. En el
\cref{chap:algebras-cuadraticas-orden-nueve} se construye
\[
\mathbb F_9=\mathbb F_3[\iota]/(\iota^2+1),
\qquad \iota^2=-1,
\]
y se demuestra que la matriz Witt realiza esa estructura cuadrática. Se trata
de una extensión algebraica finita exacta, no de una categoría visible de
expansiones reales.

\section{Carta de Hensel no filtrada, supervivencia y frontera global}

Los siguientes controles son falsadores de los resultados finitos
precedentes:
\begin{enumerate}
\item que dos estados distintos módulo \(3^T\) emitan la misma palabra de
\(T\) bloques, o que exista una palabra sin preimagen;
\item que falle la inversa explícita de
\cref{eq:inversa-hensel-finita};
\item que aparezca un racional cuya expansión canónica no sea finalmente
periódica, o un irracional cuya expansión sí lo sea;
\item que un emisor determinista finito produzca una cinta no finalmente
periódica;
\item que los censos del catálogo no sumen \(468\), \(243\) y \(104\,976\),
o que las medidas de las tres palabras no coincidan con la enumeración.
\end{enumerate}

La carta Hensel bruta
\[
\mathbb Z_3^6\longrightarrow[0,1]
\]
es continua y sobreyectiva; al añadir la parte entera, su atlas bruto cubre
\(\mathbb R\). Este resultado exacto no implica que toda fibra bruta contenga
un representante que sobreviva a los filtros adicionales. Las ventanas
finitas publicadas y los lectores nombrados de \(\pi,e,\varphi\) se verifican
dentro de HMT, relativamente a sus semillas, regiones,
orientaciones y reglas enriquecidas.

\ifdefined\HMTInternalTraceability
Para un prefijo filtrado, sea \(T_u\) el árbol de prolongaciones compatibles.
Si \(T_u\) es finitamente ramificado y posee un nodo en toda profundidad, el
lema de König produce una rama infinita; si, además, cada nivel estable es
unitario, la rama es única. Este enunciado es condicional: no se ha demostrado
la no vacuidad universal de \(T_u\) para todo prefijo. Permanecen, por tanto,
como \textsc{programa abierto} la cobertura de todas las fibras por
supervivientes filtrados y el descenso bien definido de la suma y el producto
nativos al cociente arquimediano.
\fi




\label{chap:taxonomia-monodromia}

La representación posicional de un número registra una sucesión de símbolos,
pero no conserva necesariamente el estado que determina su continuación. La
Holografía Modular Triádica distingue ambos niveles. La fase nonádica indica
la posición de una lectura dentro del ciclo de nueve pasos; el estado
enriquecido conserva, además, orientación, cociente euclídeo, transporte de
arrastre y condición de frontera. El retorno de la fase no implica, por
tanto, el retorno del estado.

Sea \(S\) un espacio de estados enriquecidos, sea
\(T:S\to S\) una transición y sea
\[
g:S\longrightarrow\mathbb Z/9\mathbb Z
\]
la fase, con \(g(Ts)=g(s)+1\). Para una base \(b\geq2\), un lector
\(d:S\to\{0,\ldots,b-1\}\) produce
\[
x(s)=\sum_{n\geq0}\frac{d(T^ns)}{b^{n+1}}.
\]
Cuando un valor admite las dos expansiones posicionales usuales, se adopta la
representación canónica que no termina en una cola constante \(b-1\). Esta
convención separa una ambigüedad de escritura de la multiplicidad propiamente
HMT, que procede de secciones enriquecidas distintas con una misma
evaluación.
El operador de retorno nonádico es la restricción
\[
M=T^9:g^{-1}(r)\longrightarrow g^{-1}(r).
\]
La diferencia entre \(T^9s=s\) y \(g(T^9s)=g(s)\) formaliza la diferencia
entre periodicidad y monodromía.

La transición del estado enriquecido de TPK puede estar dada por una relación y no por una
función. En tal caso, la notación anterior se aplica al desplazamiento sobre
el espacio de historias infinitas o a una realización determinista cuyo estado
se haya ampliado con la información necesaria para fijar la continuación; no se
presupone que un bloque visible posea un sucesor único.

\begin{definition}[Número HMT y presentación aritmética]
Un número HMT es únicamente una sección compatible
\(\boldsymbol x\) del sistema inverso de estados supervivientes. Elegido un
lector posicional \(L\) cuyo dominio contiene \(\boldsymbol x\), el par
\((\boldsymbol x,L)\) es una presentación o medición del número. El valor
arquimediano publicado por esa presentación es la intersección puntual de los
cilindros asociados cuando sus diámetros tienden a cero. Dos secciones con el
mismo valor no se identifican antes de aplicar el cociente arquimediano.
\end{definition}

\begin{definition}[Periodicidades visible y enriquecida]
La salida es eventualmente periódica cuando existen \(N\geq0\) y \(p>0\)
tales que
\[
d(T^{n+p}s)=d(T^ns),\qquad n\geq N.
\]
El estado es eventualmente periódico cuando
\[
T^{N+p}s=T^Ns.
\]
La segunda condición implica la primera. La recíproca exige que el lector
separe las colas de la órbita considerada, es decir, que
\[
 \bigl(d(T^{n+k}s)\bigr)_{k\geq0}
 =\bigl(d(T^{m+k}s)\bigr)_{k\geq0}
 \quad\Longrightarrow\quad T^ns=T^ms.
\]
\end{definition}

\begin{theorem}[Clasificación por retorno de fase]
\label{thm:clasificacion-retorno-fase}
Sea \(s\in S\).
\begin{enumerate}
  \item Si la órbita de \(s\) alcanza un estado cuya salida futura es
  constantemente nula, entonces \(x(s)\) posee una expansión terminante y es
  racional.
  \item Si el estado es eventualmente periódico, entonces la salida es
  eventualmente periódica y \(x(s)\in\mathbb Q\).
  \item Si el lector separa las colas de la órbita y la órbita de \(s\) no es
  eventualmente periódica, entonces la salida no es eventualmente periódica
  y \(x(s)\notin\mathbb Q\).
  \item La no periodicidad no distingue entre irracionalidad algebraica y
  trascendencia. Esta distinción requiere un certificado polinómico o un
  teorema de trascendencia adicional.
\end{enumerate}
\end{theorem}

\begin{proof}
Los dos primeros apartados son la caracterización clásica de los racionales
mediante expansiones posicionales terminantes o eventualmente periódicas. En
el tercero, una eventual periodicidad de la salida produciría dos colas
iguales. La hipótesis de separación identificaría entonces los estados
correspondientes, contradiciendo la no periodicidad de la órbita. El cuarto
apartado sigue de la existencia de irracionales algebraicos y trascendentes
con expansiones no eventualmente periódicas.
\end{proof}

\begin{corollary}[Monodromía nonádica y aperiodicidad]
Supóngase que la fase retorna cada nueve pasos, que la órbita de retorno
\[
 \{M^ks:k\geq0\},\qquad M=T^9,
\]
no es eventualmente periódica y que el lector separa las colas. Entonces la
salida no es eventualmente periódica y el número producido es irracional. La
periodicidad de la fase convive en este caso con una evolución aperiódica del
estado enriquecido.
\end{corollary}

La hipótesis no puede sustituirse por una sola desigualdad \(M(s)\ne s\), ni
por un número finito de retornos distintos. Relativamente al estado inicial,
la transición y el lector declarados, la ventana vigente certifica nueve
retornos de fase con cambio de estado; demuestra monodromía finita, pero no
por sí sola la aperiodicidad de toda la órbita.

\ifdefined\HMTRepetirResultadosTaxonomicos
\section{El cuadrado local de \texorpdfstring{\(e\)}{e} y la ruptura de memoria}

Para una palabra decimal \(a_1\cdots a_n\), llamaremos \emph{cuadrado local}
de longitud \(\ell\) en la posición \(i\) a una factorización
\[
 a_i\cdots a_{i+2\ell-1}=vv,
 \qquad |v|=\ell.
\]
Esta noción pertenece a combinatoria de palabras y no debe confundirse con
un período de la cola infinita.

\begin{theorem}[Caracterización finita del cuadrado \(1828^2\)]
Considérense las \(243\) palabras del catálogo N33 y el lector congelado
\(w_6\mapsto w_{30}\mapsto d_1d_2d_3d_4\), con cuatro tríadas en base mil.
Existe una única palabra cuya lectura contiene un cuadrado de longitud cuatro
iniciado en la segunda cifra decimal. Es
\[
 w_e=201101,
 \qquad
 718281828459=7\,\underbrace{1828\,1828}_{(1828)^2}\,459.
\]
La repetición no se prolonga a una tercera copia: la cifra siguiente requerida
sería \(1\), mientras que la cifra efectiva es \(4\). En todo el catálogo sólo
existe otro cuadrado de longitud cuatro,
\[
 575157518472=\underbrace{5751\,5751}_{\text{cuadrado de palabra}}\,8472,
\]
pero comienza en la primera cifra.
\end{theorem}

\begin{proof}
La enumeración exhaustiva evalúa las \(243\) palabras mediante aritmética
entera exacta. El perfil de cuadrados para longitudes \(1,\ldots,6\) contiene,
respectivamente,
\[
240, 21, 3, 2, 0, 0
\]
ocurrencias, distribuidas entre
\[
153, 19, 3, 2, 0, 0
\]
palabras. Las dos ocurrencias de longitud cuatro son exactamente las escritas
en el enunciado. Dos enumeraciones enteras independientes reproducen el mismo
censo y sus salidas completas coinciden término a término.
\end{proof}

La palabra \(w_e\) posee además la cadena de elevación
\[
201101\mid121221\mid102011\mid012222\mid102011.
\]
El tercer y el quinto bloque visibles coinciden. No coinciden, sin embargo,
los estados que los transportan. Antes de ambas emisiones, sus reducciones
módulo \(27\) son
\[
(25,11,7,17,8,16),
\qquad
(7,8,4,23,26,7),
\]
y los cocientes de Hensel posteriores son
\[
(6,13,9,2,13,1),
\qquad
(15,0,3,19,23,25).
\]
Por tanto, el retorno \(102011\) es un retorno de la proyección visible, no
del estado enriquecido. Esta es la formulación exacta de la intuición de una
repetición local que se rompe: el cociente recuerda una genealogía que el
bloque residual no conserva.

El resultado pertenece al lector enriquecido ya fijado por el sistema
APP--TRIT--TPK\@. El censo finito caracteriza exhaustivamente la aparición del
cuadrado local; la prolongación coinductiva de la sección pertenece a la
conexión nonádica construida en el
\cref{chap:numero-heisenberg-d108}. La propiedad combinatoria del prefijo y
la generación de la expansión completa son, por tanto, dos niveles
compatibles de una misma genealogía.

Esta formulación proporciona el marco exacto para interpretar el comienzo
decimal de \(e\),
\[
e=2.718281828459\ldots.
\]
La expresión correcta ya no es ``semiperiodicidad'', sino \emph{cuadrado local
de palabra con retorno visible y ruptura de memoria}.

\section{2 medidas finitas del emisor \(468\to243\): frecuencia y prioridad estructural}

No existe una distribución uniforme sobre todos los números reales. Sí
existe una medida canónica sobre el catálogo finito cuando se toma uniforme
el conjunto de \(104\,976\) semillas TPK\@. Para una palabra visible \(w\), se
define
\[
\mu_{\mathrm{TPK}}(w)
=\frac{\#\{\text{semillas que proyectan sobre }w\}}{104\,976}.
\]
La medida no es uniforme. Las \(243\) palabras poseen pesos comprendidos
entre \(144\) y \(2088\), con nueve valores de multiplicidad distintos.
El histograma exacto, escribiendo el peso como base y el número de palabras
como exponente, es
\[
144^{120},\ 288^{54},\ 432^{18},\ 576^9,\ 864^{15},\
1008^6,\ 1440^3,\ 1944^{12},\ 2088^6.
\]

Para las tres palabras distinguidas se obtiene
\[
\mu_{\mathrm{TPK}}(010211)=\frac{1008}{104\,976}=\frac7{729},
\]
\[
\mu_{\mathrm{TPK}}(201101)
=\mu_{\mathrm{TPK}}(121200)
=\frac{144}{104\,976}=\frac1{729}.
\]
La palabra asociada a \(\pi\) tiene cinco preimágenes \(U_6\) y peso siete
veces mayor que las palabras asociadas a \(e\) y \(\varphi\). Este resultado
demuestra una jerarquía combinatoria en el catálogo; no identifica por sí
solo las palabras con las constantes clásicas ni establece una probabilidad
sobre \(\mathbb R\).

Tampoco define por sí solo una prioridad total: seis palabras tienen peso
\(1008\), mientras que \(e\) y \(\varphi\) pertenecen a la clase mínima de
peso \(144\), compartida por otras \(118\) palabras. La distinción HMT debe
expresarse mediante un perfil estructural, por ejemplo
\[
\mathcal P(w)=
\bigl(m_{\rm semilla}(w),m_{U_6}(w),r_{\rm supervivencia}(w),
s_{\rm frontera}(w),s_{\rm excepcional}(w)\bigr),
\]
con cada coordenada definida antes de aplicar etiquetas externas. Un escalar
de ``prioridad'' exige además fijar, también de antemano, cómo se ordenan o
ponderan esas coordenadas.
\fi

\section{Clasificación de las constantes por antecedente genealógico, operación generadora y codominio}

La clasificación HMT separa propiedades que la representación ordinaria
reúne bajo un solo valor:
\begin{center}
\begin{tabular}{p{0.25\textwidth}p{0.64\textwidth}}
\toprule
Clase & Condición estructural \\
\midrule
Terminante & La salida alcanza una cola nula. \\
Racional periódico & La salida es eventualmente periódica. \\
Racional monodrómico & La salida es periódica, pero el estado enriquecido
recorre una fibra no trivial que el lector visible olvida. \\
Irracional coinductivo & La salida no es eventualmente periódica y los
cilindros determinan un único valor. \\
Recurrente local & Una palabra finita contiene retornos o cuadrados, sin que
ello imponga periodicidad a la cola. \\
Algebraico & El valor satisface un polinomio no nulo con coeficientes enteros;
la ruta debe acompañarse de ese certificado. \\
Trascendente & El valor no satisface ningún polinomio entero no nulo; la
clasificación exige un teorema de trascendencia. \\
Distinguido HMT & La sección satisface un predicado estructural congelado de
fibra, frontera, supervivencia o simetría; es una clase transversal a las
anteriores. \\
\bottomrule
\end{tabular}
\end{center}

Así, HMT no altera las definiciones aritméticas correctas de racionalidad,
algebraicidad y trascendencia. Las refina añadiendo la genealogía del estado,
la monodromía de fase y la multiplicidad de las fibras que la evaluación
arquimediana había olvidado.

En particular, \(\varphi\) es algebraico porque satisface
\(X^2-X-1=0\), mientras que la trascendencia de \(e\) y \(\pi\) procede de
teoremas externos clásicos. Estas propiedades son compatibles con que las tres
secciones sean distinguidas por HMT\@. El cierre dodecafásico construye más
adelante la proyección racional exacta
\(\alpha_{12}=\pi_{12}(\alpha_{\mathrm{HMT}})\), y la prolongación compatible
construye la sección coinductiva \(\alpha_{\mathrm{HMT}}\), en el
\cref{chap:alpha-rev7}; ninguna de ellas se usa en esta taxonomía como entrada
ni se publica aquí como una segunda residencia. El núcleo \(E_{21/22}\)
produce una raíz analítica distinta, cuya concordancia con la vía entera se
demuestra en el codominio de \(\operatorname{Tr}_9\). Sección, proyección
racional, raíz analítica, lector y valor conservan así sus tipos propios.

Finalmente, una ``dirección infinitesimal'' requiere tipado propio. Para un
lector finito \(r_m\), dos secciones distintas en una misma fibra de \(r_m\)
definen una diferencia invisible a profundidad \(m\), aunque activa en una
prolongación posterior. No constituyen por ello un real infinitesimal no nulo.
Si dos secciones coinciden bajo todas las truncaciones fieles del sistema
inverso, son la misma sección. Un infinitesimal aritmético genuino exigiría
construir una extensión ordenada o valuada no arquimediana y demostrar que las
operaciones HMT descienden a ella.



\section[N\'umero como secci\'on compatible]{N\'umero como secci\'on
compatible: clasificaci\'on aritm\'etica, evaluaci\'on arquimediana y
genealog\'ia}
\label{p3-83-c44-s1:sec:ontologia-numeros-rev11}
\index{n\'umero!valor y genealog\'ia}
\index{taxonom\'ia num\'erica}

La clasificaci\'on num\'erica de HMT conserva las especies aritm\'eticas
ordinarias y a\~nade una coordenada que la evaluaci\'on posicional descarta: la
genealog\'ia de la secci\'on que publica el valor. El resultado es una
clasificaci\'on conjunta. El primer eje responde qu\'e n\'umero se obtiene; el
segundo registra mediante qu\'e sucesi\'on compatible de estados, retornos,
acarreos y cierres se obtiene. Esta distinci\'on permite formular en un mismo
lenguaje los enteros, las expansiones racionales, los irracionales, la unidad
imaginaria, las regularizaciones y las constantes espectrales sin identificar
objetos que s\'olo comparten un escalar.

\subsection{La evaluaci\'on arquimediana separa la secci\'on de su valor}

Sea
\[
 X_{\infty}^{\rm surv}
 =\varprojlim_n X_n^{\rm surv}
\]
un sistema inverso de estados supervivientes con sus aplicaciones de
restricci\'on. Una secci\'on compatible es una familia
\(\boldsymbol x=(x_n)_{n\geq0}\) que satisface
\(\rho_{n+1,n}(x_{n+1})=x_n\) a toda profundidad. Un lector posicional
\(L\) asigna a cada prefijo un cilindro \(I_n(\boldsymbol x;L)\). Cuando los
cilindros son compatibles, sus di\'ametros convergen a cero y se fija la
convenci\'on can\'onica de extremos, se define
\begin{equation}
 \operatorname{Val}_{L}(\boldsymbol x)
 :=\text{el punto de }
 \bigcap_{n\geq0}\overline{I_n(\boldsymbol x;L)}.
 \label{p3-83-c44-s1:eq:valor-lector-seccion-rev11}
\end{equation}

\HMTClaim{HMT-R-NUMBER-SECTION-VALUE}
\begin{definition}[N\'umero, presentaci\'on y valor en HMT]
\label{p3-83-c44-s1:def:numero-presentacion-valor-rev11}
Un \emph{n\'umero HMT} es una secci\'on compatible
\(\boldsymbol x\in X_{\infty}^{\rm surv}\). El par
\((\boldsymbol x,L)\) es una \emph{presentaci\'on} del n\'umero y
\(\operatorname{Val}_{L}(\boldsymbol x)\) es el valor publicado por esa
presentaci\'on. La fibra
\[
 \operatorname{Val}_{L}^{-1}(r)
\]
conserva las genealog\'ias que el cociente arquimediano representa mediante
el mismo valor \(r\).
\end{definition}

La capacidad de la carta y la selecci\'on de una secci\'on son resultados
tipados distintos. El \cref{thm:cobertura-arquimediana-bruta} construye una
aplicaci\'on continua y sobreyectiva
\[
 P_{\rm raw}:\mathbb Z_3^6\longrightarrow[0,1],
\]
y el \cref{cor:atlas-hmt-reales} la globaliza a
\[
 \widehat P_{\rm raw}:\mathbb Z\times\mathbb Z_3^6
 \longrightarrow\mathbb R.
\]
Estos mapas prueban la capacidad representacional completa de la carta bruta.
El predicado de supervivencia y los lectores que distinguen
\(\pi,e,\varphi\) seleccionan, dentro de esa capacidad, las genealog\'ias
estructurales correspondientes.

\subsection{Terminaci\'on y periodicidad determinan la racionalidad}

Para la expansi\'on can\'onica en base mil
\[
 x=\sum_{n\geq1}\frac{b_n}{1000^n},
 \qquad b_n\in\{0,\ldots,999\},
\]
se adopta la convenci\'on que elimina la duplicidad producida por una cola
finalmente constante de bloques \(999\). El
\cref{thm:taxonomia-base-mil} proporciona entonces la clasificaci\'on exacta
siguiente.

\HMTClaim{HMT-R-NUMBER-DECIMAL-CLASSIFICATION}
\begin{theorem}[Clasificaci\'on visible en la carta decimal c\'ubica]
\label{p3-83-c44-s1:thm:clasificacion-visible-rev11}
Sea \(x=a/q\in[0,1)\) una fracci\'on irreducible cuando \(x\) sea racional.
Entonces:
\begin{enumerate}
\item la expansi\'on termina si y s\'olo si \(q\mid1000^N\) para alg\'un
      \(N\), equivalencia que restringe los factores primos de \(q\) a
      \(2\) y \(5\);
\item la expansi\'on es finalmente peri\'odica si y s\'olo si
      \(x\in\mathbb Q\);
\item la expansi\'on es no finalmente peri\'odica si y s\'olo si
      \(x\in\mathbb R\setminus\mathbb Q\).
\end{enumerate}
La algebraicidad y la trascendencia constituyen un segundo eje: un real es
algebraico si anula un polinomio no nulo de \(\mathbb Q[X]\), y es
trascendente si no anula ninguno.
\end{theorem}

\begin{proof}
Una expansi\'on de longitud \(N\) tiene denominador divisor de \(1000^N\), y
la rec\'iproca se obtiene expresando \(a/q\) con ese denominador. Para un
racional, la divisi\'on sucesiva posee un conjunto finito de restos: un resto
nulo produce terminaci\'on y la repetici\'on de un resto produce periodicidad.
Una cola peri\'odica es, rec\'iprocamente, una serie geom\'etrica racional. La
clasificaci\'on polin\'omica es independiente de la base posicional.
\end{proof}

La din\'amica enriquecida distingue adem\'as periodicidad de salida y
periodicidad de estado. Si \(T\) transporta el estado y \(d\) publica un
bloque, las condiciones
\[
 d(T^{n+p}s)=d(T^n s),
 \qquad
 T^{n+p}s=T^n s
\]
definen, respectivamente, retorno visible y retorno del estado. La segunda
implica la primera. La implicaci\'on inversa requiere que el lector separe las
colas de la órbita, conforme al
\cref{thm:clasificacion-retorno-fase}. Por ello, una fase non\'adica peri\'odica
puede transportar una memoria no peri\'odica.

\begin{longtable}{@{}>{\raggedright\arraybackslash}p{.20\textwidth}
                    >{\raggedright\arraybackslash}p{.30\textwidth}
                    >{\raggedright\arraybackslash}p{.42\textwidth}@{}}
\caption{Taxonom\'ia conjunta del valor publicado y del estado enriquecido.}
\label{p3-83-c44-s1:tab:taxonomia-conjunta-numeros-rev11}\\
\toprule
\textbf{Clase} & \textbf{Condici\'on del valor} &
\textbf{Condici\'on geneal\'ogica o demostraci\'on requerida}\\
\midrule
\endfirsthead
\toprule
\textbf{Clase} & \textbf{Condici\'on del valor} &
\textbf{Condici\'on geneal\'ogica o demostraci\'on requerida}\\
\midrule
\endhead
Entero & Elemento de \(\mathbb Z\). & Palabra ternaria balanceada finita y
\'unica; signo y magnitud permanecen tipados.\\
Decimal finito & Racional cuyo denominador reducido divide alguna potencia
de \(1000\). & La salida alcanza una cola nula.\\
Racional peri\'odico & Expansi\'on finalmente peri\'odica. & Un estado
eventualmente peri\'odico basta; una salida peri\'odica puede conservar una
fibra interna adicional.\\
Racional monodr\'omico & Mismo valor racional que la fila anterior. & La
salida retorna mientras el estado enriquecido recorre una fibra que el lector
visible olvida.\\
Irracional coinductivo & Expansi\'on no finalmente peri\'odica. & Una familia
infinita de cilindros compatibles fija un único valor; la aperiodicidad
global exige una prueba sobre toda la órbita, no una ventana finita.\\
Algebraico & Ra\'iz de un polinomio no nulo de \(\mathbb Q[X]\). & Se adjunta
el polinomio y se comprueba su anulaci\'on.\\
Trascendente & Real exterior a la clausura algebraica de \(\mathbb Q\). & Se
adjunta un teorema de trascendencia; la no periodicidad decimal por s\'i sola
s\'olo demuestra irracionalidad.\\
Distinguido HMT & Cualquiera de las clases aritm\'eticas anteriores. & La
secci\'on satisface un predicado congelado de fibra, frontera, supervivencia,
simetr\'ia o cierre.\\
\bottomrule
\end{longtable}

Esta tabla sit\'ua en ejes diferentes dos preguntas que deben contestarse
separadamente. La clasificaci\'on ordinaria determina las propiedades del valor;
el lector HMT determina la cadena estructural que genera la secci\'on. Por
ejemplo,
\[
 \varphi^2-\varphi-1=0
\]
clasifica \(\varphi\) como irracional algebraico, mientras que la
trascendencia de \(e\) y de \(\pi\) pertenece a teoremas cl\'asicos externos.
Las tres pueden, simult\'aneamente, ser secciones distinguidas por predicados
HMT diferentes. Sus truncaciones impresas de mil o diez mil cifras son
racionales finitos que demuestran los prefijos calculados de una salida prolongable; no
reemplazan el objeto infinito que representan.

El cierre dodecaf\'asico define, por su parte, la proyecci\'on racional finita
\begin{equation}
 \alpha_{12}=\pi_{12}(\alpha_{\rm HMT})
 =\frac{7297352569283800997285105472380663}{10^{36}}.
 \label{p3-83-c44-s1:eq:alpha-racional-taxonomia-rev11}
\end{equation}
Esta proyecci\'on es exacta y racional. El objeto
\(\alpha_{\rm HMT}=\varprojlim_NA^{(N)}\) es la secci\'on infinita construida
por la recurrencia compatible; sus truncaciones racionales no la sustituyen.
La derivaci\'on estructural del l\'imite y la clasificaci\'on de cada
proyecci\'on finita son afirmaciones complementarias.

\subsection{La carta ternaria balanceada representa los enteros
biyectivamente}

Sea
\[
 \mathscr W_{\rm bal}
 =\{\varnothing\}\cup
 \{a_1\cdots a_m:
 a_j\in\{-1,0,+1\},\ a_1\neq0\}
\]
y def\'inase
\[
 b(\varnothing)=0,
 \qquad
 b(wa)=3b(w)+a.
\]
El \cref{pf84:C03} demuestra que
\begin{equation}
 b:\mathscr W_{\rm bal}\xrightarrow{\ \sim\ }\mathbb Z
 \label{p3-83-c44-s1:eq:carta-entera-balanceada-rev11}
\end{equation}
es una biyecci\'on. El último trit es el representante único de la
clase m\'odulo tres y el cociente \((n-a)/3\) reduce estrictamente el valor
absoluto hasta alcanzar cero. El primer trit no nulo fija el signo. Esta carta
es una representaci\'on aritm\'etica exacta; la pertenencia de una palabra a una
secci\'on superviviente incorpora adem\'as las condiciones de ruta,
orientaci\'on y frontera.

Para \(n>0\), las coordenadas
\[
 \rho_9(n)=1+((n-1)\bmod9),
 \qquad
 \kappa_9(n)=\frac{n-\rho_9(n)}9
\]
realizan la descomposici\'on \(n=9\kappa_9(n)+\rho_9(n)\). La fase
\(\rho_9\in\{1,\ldots,9\}\) publica la clase residual cero como \(9\), y
\(\kappa_9\) conserva el n\'umero de vueltas. Con el signo
\(\varepsilon\in\{-1,+1\}\),
\[
 n\longleftrightarrow
 \bigl(\varepsilon,\kappa_9(|n|),\rho_9(|n|)\bigr)
\]
es biyectivo sobre \(\mathbb Z\setminus\{0\}\); el cero permanece como punto
base. El retorno visible \(9\to1\) es as\'i una transici\'on de fase con
incremento de memoria, no una identificaci\'on de enteros consecutivos.

\subsection{Estructura compleja finita inducida por el endomorfismo de Witt \texorpdfstring{\(J^2=-I_6\)}{J²=-I₆}}

La unidad imaginaria ordinaria es el elemento
\(i\in\mathbb C\) que satisface \(i^2=-1\); es algebraico de grado dos sobre
\(\mathbb R\) y sobre \(\mathbb Q\). La realizaci\'on HMT pertenece a un
dominio finito diferente. Si \(C_P\) es la matriz de conferencia de Paley y
\(J=A_W\) su reducci\'on m\'odulo tres en
\(\operatorname{End}_{\mathbb F_3}(\mathbb F_3^6)\), entonces
\begin{equation}
 \boxed{J^2=2I_6=-I_6.}
 \label{p3-83-c44-s1:eq:raiz-interna-menos-uno-ontologia-rev11}
\end{equation}
Como \(X^2+1\) es irreducible en \(\mathbb F_3[X]\), la extensi\'on
\[
 \mathbb F_9=\mathbb F_3[\iota]/(\iota^2+1)
\]
separa los espacios propios de \(J\) asociados a \(\pm\iota\). El resultado,
demostrado en \cref{sec:raiz-interna-menos-uno-rev6}, dota al m\'odulo ternario de una
estructura compleja finita y organiza dos orientaciones conjugadas. Su fuerza
es \emph{exacta interna, demostrada con la matriz Paley--Witt declarada}; la
comparaci\'on con \(\mathbb C\) conserva la distinci\'on entre ambos cuerpos.

\subsection{\(6\) operadores realizan el racional
\texorpdfstring{\( -1/12 \)}{-1/12} en dominios distintos}

La regularizaci\'on zeta asigna a la sucesi\'on de enteros positivos la parte
finita obtenida por continuaci\'on meromorfa:
\begin{equation}
 \sum_{n\geq1}^{[\zeta]}n
 :=\zeta(-1)=-\frac1{12}.
 \label{p3-83-c44-s1:eq:regularizacion-zeta-menos-doce-rev11}
\end{equation}
El super\'indice \([\zeta]\) forma parte del tipo de la expresi\'on. La suma
parcial ordinaria satisface
\(\sum_{n=1}^{N}n=N(N+1)/2\) y crece sin cota; por ello
\eqref{p3-83-c44-s1:eq:regularizacion-zeta-menos-doce-rev11} es un valor regularizado, no
el l\'imite de esas sumas parciales.

HMT contiene adem\'as seis realizaciones algebraicas exactas del mismo
racional. La coincidencia escalar permite compararlas; cada objeto conserva su
dominio, su mapa y su normalizaci\'on.

\begin{longtable}{@{}>{\raggedright\arraybackslash}p{.23\textwidth}
                    >{\raggedright\arraybackslash}p{.43\textwidth}
                    >{\raggedright\arraybackslash}p{.26\textwidth}@{}}
\caption{Realizaciones tipadas de \(-1/12\).}
\label{p3-83-c44-s1:tab:realizaciones-menos-doce-rev11}\\
\toprule
\textbf{Dominio} & \textbf{Mapa y evaluaci\'on} &
\textbf{Conclusi\'on y normalizaci\'on}\\
\midrule
\endfirsthead
\toprule
\textbf{Dominio} & \textbf{Mapa y evaluaci\'on} &
\textbf{Conclusi\'on y normalizaci\'on}\\
\midrule
\endhead
Incidencia marcada hexada--octada &
\(F_6\hookrightarrow F_8\), con
\(\lvert F_6\rvert=90\), \(\lvert F_8\rvert=120\), y
\[
 (90-120)/(24\binom62)=-1/12.
\]
& Defecto combinatorio exacto para la normalizaci\'on declarada.\\
Periodicidades dodecaf\'asicas &
Un paso de \(30^\circ\) produce
\[
 \begin{aligned}
 \frac{30}{120}-\frac{30}{90}
 &=\frac14-\frac13,\\
 &=-\frac1{12}.
 \end{aligned}
\]
& Diferencia orientada exacta de periodos.\\
Operador de Gram &
Para \(K_{60,81}=\operatorname{diag}(0,12)\), la pseudoinversa sobre el
rango no nulo vale \(1/12\); la orientaci\'on negativa publica \(-1/12\). &
Identidad operatorial exacta sobre la componente no nula.\\
Segunda diferencia de escala &
El c\'alculo produce \(C(L)=8/81\) y el extractor declarado
\(R(L)=-(27/32)C(L)=-1/12\). & Exacto despu\'es de fijar el factor de
normalizaci\'on \(-27/32\); dicho factor es una primitiva expl\'icita.\\
Cociente modular eta &
\[
 \begin{aligned}
 \frac{\eta(\tau)}{\eta(3\tau)}
 &=q^{-1/12}\\[-1mm]
 &\quad{}\times
 \prod_{n\geq1}(1+q^n+q^{2n})^{-1}.
 \end{aligned}
\]
& Exponente modular exacto, procedente de \(1/24-3/24\).\\
Proyectores del ciclo de doce fases &
Para \(D_r=I-S^r\),
\[
 \begin{aligned}
 \tau(\Pi_{\ker D_3}-\Pi_{\ker D_4})
 &=\frac{3-4}{12},\\
 &=-\frac1{12}.
 \end{aligned}
\]
& Defecto transversal exacto entre los pasos tres y cuatro.\\
\bottomrule
\end{longtable}

Las seis realizaciones se demuestran en
\cref{sec:menos-doce-realizaciones-rev6}. La regularizaci\'on zeta, el defecto
de incidencia, la pseudoinversa, el extractor de escala, el cociente eta y la
traza de proyectores son aplicaciones diferentes. Una identificaci\'on de sus
dominios requerir\'ia un entrelazador adicional. Asimismo, el t\'ermino
\(+1/12\) de la f\'ormula de Glaisher--Kinkelin pertenece a la regularizaci\'on
del hiperfactorial y conserva ese tipo propio.

\subsection{Descomposición multiplicativa de los primos como modos irreducibles del TPK}
\index{primo!irreducible}
\index{reloj primo}

Un primo ordinario es un entero \(p>1\) cuyas factorizaciones
\(p=ab\), con \(a,b\in\mathbb N\), contienen necesariamente una unidad. El
coproducto reducido
\[
 \Delta_{\times}^{\rm red}e_n
 =\sum_{\substack{ab=n\\a,b>1}}e_a\otimes e_b
\]
convierte esa definici\'on en un selector espectral. Seg\'un el
\cref{pf84:C07}, el operador positivo
\[
 N_{\times}^{\rm red}
 =(\overline{\Delta_{\times}^{\rm red}})^*
   \overline{\Delta_{\times}^{\rm red}}
\]
satisface
\(N_{\times}^{\rm red}e_n=d_{\rm red}(n)e_n\), y su n\'ucleo, despu\'es de
retirar \(e_1\), selecciona exactamente los primos. El entrelazador
multiplicativo transporta el mismo selector al producto nativo de APP\@. Este
selector es la proyección espectral sobre
\(\ker N_{\times}^{\rm red}\ominus\mathbb C e_1\), que coincide exactamente
con el subespacio generado por los primos y no depende de reconocer listas
decimales.

La factorizaci\'on única proporciona adem\'as la realizaci\'on de Fock
\[
 U_F:\Gamma_s(\ell^2(\mathbb P))
 \xrightarrow{\ \sim\ }\ell^2(\mathbb N_{\geq1}),
 \qquad
 U_F\,d\Gamma(H_{\mathbb P})\,U_F^*=\log Q,
\]
donde \(H_{\mathbb P}e_p=(\log p)e_p\). Cada primo es un modo
multiplicativamente irreducible; un entero compuesto es una ocupaci\'on finita
de esos modos.

Para \(n\) coprimo con \(10\), su reloj decimal es el orden multiplicativo
\[
 \tau_n=\operatorname{ord}_n(10).
\]
Si \(p\) es primo y \(p\notin\{2,3,5\}\), el
\cref{thm:cierre-nonadico-reloj-primo-rev6} demuestra
\begin{equation}
 \tau_p\mid p-1,
 \qquad
 M_p:=\frac{10^{\tau_p}-1}{p}\in\mathbb N,
 \qquad
 M_p\equiv0\pmod9.
 \label{p3-83-c44-s1:eq:reloj-primo-cierre-nonadico-rev11}
\end{equation}
Para un compuesto,
\[
 \operatorname{ord}_{n}(10)
 =\operatorname{lcm}_{p^a\parallel n}\operatorname{ord}_{p^a}(10).
\]
Por tanto, el reloj compuesto sincroniza relojes de potencias primas, mientras
que el primo aporta un componente irreducible. La congruencia non\'adica
\eqref{p3-83-c44-s1:eq:reloj-primo-cierre-nonadico-rev11} describe el cierre de fase; el
selector de primalidad reside en el coproducto reducido. La separaci\'on de
estas dos funciones evita atribuir capacidad selectiva al periodo decimal por
s\'i solo.

\subsection{Cada constante queda asociada a su operador de extracci\'on}

Los ciclos de orden \(3^m\), el l\'imite de sus trazas de calor y la
transformada theta--Mellin construyen el funcional
\begin{equation}
 \mathcal Z_3(s)
 =\frac{\displaystyle
 \int_0^\infty\Theta_+(x)x^{s/2-1}\,\mathrm dx}
 {\pi^{-s/2}\Gamma(s/2)}
 =\zeta(s),
 \qquad \Re s>1.
 \label{p3-83-c44-s1:eq:funcional-triadico-resumen-rev11}
\end{equation}
La inversi\'on theta prolonga merom\'orficamente el funcional. La partici\'on
residual
\[
 \zeta(s)=Z_0(s)+Z_1(s)+Z_2(s)
\]
separa agregado, escala y orientaci\'on. Una sola fuente operatorial admite
varios mapas de extracci\'on ---parte finita, coeficiente de Laurent,
evaluaci\'on, derivada regularizada o car\'acter---; la coincidencia de fuente
no identifica las constantes resultantes.

\newpage
\enlargethispage{2\baselineskip}
\begingroup
\small
\fontsize{10pt}{11.2pt}\selectfont
\setlength{\LTpre}{0.5\baselineskip}
\setlength{\LTpost}{0.5\baselineskip}
\begin{longtable}{@{}>{\raggedright\arraybackslash}p{.15\textwidth}
                    >{\raggedright\arraybackslash}p{.245\textwidth}
                    >{\raggedright\arraybackslash}p{.305\textwidth}
                    >{\raggedright\arraybackslash}p{.195\textwidth}@{}}
\caption{Definici\'on, mapa de extracci\'on y alcance matem\'atico de las
constantes num\'ericas complementarias.}
\label{p3-83-c44-s1:tab:constantes-operadores-ontologia-rev11}\\
\toprule
\textbf{Objeto} & \textbf{Definici\'on est\'andar} &
\textbf{Mapa y referencia matem\'atica} & \textbf{Conclusi\'on y condici\'on}\\
\midrule
\endfirsthead
\toprule
\textbf{Objeto} & \textbf{Definici\'on est\'andar} &
\textbf{Mapa y referencia matem\'atica} & \textbf{Conclusi\'on y condici\'on}\\
\midrule
\endhead
Euler--Mascheroni \(\gamma\) &
Para \(H_N=\sum_{n\leq N}n^{-1}\),
\(\gamma=\lim_{N\to\infty}(H_N-\log N)\), equivalente a la parte finita de
\(\zeta\) en \(s=1\). &
Si \(C_r=\operatorname*{FP}_{s=1}Z_r(s)\), entonces
\(\gamma=C_0+C_1+C_2\), por
\eqref{eq:residual-three-covectors}. &
Identidad exacta del funcional residual demostrado con el operador y la
normalizaci\'on declarados.\\
Constantes de Stieltjes \(\gamma_n\) &
Si \(\zeta(1+z)-z^{-1}=\sum_{n\geq0}a_nz^n\), se definen por
\(\gamma_n=(-1)^n n!a_n\). &
Los coeficientes residuales \(c_{r,n}\) satisfacen
\(\gamma_n=(-1)^n n!\sum_r c_{r,n}\) y
\(L^{(n)}(1,\chi_{-3})/n!=c_{1,n}-c_{2,n}\): es la descomposición
coeficiente a coeficiente de \eqref{eq:stieltjes} por las tres clases
residuales de \eqref{eq:residual-three-covectors}. &
Descomposici\'on exacta; las extracciones num\'ericas comprueban los
coeficientes, no su definici\'on.\\
Ap\'ery \(\zeta(3)\) &
Valor de la zeta de Riemann en \(s=3\). &
\(\mathcal Z_3(3)=\zeta(3)\), con prefijos acotados por intervalos racionales;
su residencia espectral dentro de la torre graduada se fija en
\cref{eq:bendersky,eq:odd-zeta-bendersky}. &
Evaluaci\'on exacta del funcional; su irracionalidad procede del teorema
cl\'asico de Ap\'ery \cite{Apery1979}.\\
Glaisher--Kinkelin \(A_{\rm GK}\) &
Constante del crecimiento regularizado del hiperfactorial. &
\(A_{\rm GK}=\exp(1/12-\mathcal Z_3'(-1))\), por
\cref{eq:bendersky,eq:odd-zeta-bendersky}. &
Derivada regularizada exacta; \(1/12\) pertenece a esta normalizaci\'on
hiperfactorial.\\
Euler--Kronecker \(\gamma_{F_{-3}}\) &
Parte finita logar\'itmica de la zeta de Dedekind del campo
\(F_{-3}=\mathbb Q(\sqrt{-3})\). &
El car\'acter residual \(\chi_{-3}\) selecciona el campo y produce
\(\gamma_{F_{-3}}=\gamma+L'(1,\chi_{-3})/L(1,\chi_{-3})\), por
\eqref{eq:euler-kronecker}. &
Identidad exacta; la elecci\'on del campo procede del car\'acter m\'odulo tres.\\
Meissel--Mertens \(B_1\) &
Para \(S_P(x)=\sum_{p\leq x}p^{-1}\),
\(B_1=\lim_{x\to\infty}(S_P(x)-\log\log x)\). &
Es la parte finita de la traza prima truncada sobre el sector de una part\'icula
del Hamiltoniano de Fock, por \eqref{eq:meissel-mertens}; el selector de
modos primos se fija antes de efectuar la regularizaci\'on. &
Representaci\'on operatorial exacta una vez fijados los modos primos; el
selector de irreducibles es anterior.\\
Feigenbaum \((\delta,\alpha_{\rm F})\) &
Constantes universales de escala del punto fijo de renormalizaci\'on para la
cascada cuadr\'atica de duplicaci\'on de periodo. &
La familia dodecaf\'asica declarada produce a nivel ocho
\(\delta^{\mathscr B}_8(0)=4.66921218915\ldots\) y
\(\alpha^{\mathscr B}_8(0)=2.50296847988\ldots\), por
\cref{sec:p3-final:cierre-feigenbaum}. &
C\'alculo finito exacto hasta \(n=8\); la convergencia al punto fijo universal
requiere demostrar la convergencia del operador de renormalizaci\'on.
\(\alpha_{\rm F}\) y \(\alpha_{\rm HMT}\) son
objetos distintos.\\
Catalan \(G\) &
\(G=L(2,\chi_{-4})=\sum_{n\geq0}(-1)^n/(2n+1)^2\). &
El lector natural es el car\'acter impar m\'odulo cuatro evaluado en \(2\),
por \cref{hmtctx:catalan:sec:combinaciones-lectores-vecinos-rev6}. &
Identidad cl\'asica exacta con operador localizado; no es una salida del
funcional residual m\'odulo tres.\\
Khinchin \(K\) &
Media geom\'etrica casi segura de los d\'igitos de fracci\'on continua para el
sistema de Gauss. &
El lector es la media erg\'odica de \(\log a_1\) para
\(T(x)=x^{-1}-\lfloor x^{-1}\rfloor\) y su medida invariante, por
\cref{hmtctx:khinchin:sec:combinaciones-lectores-vecinos-rev6}. &
Teorema erg\'odico cl\'asico incorporado; requiere una din\'amica distinta de
la torre espectral tri\'adica.\\
\bottomrule
\end{longtable}
\endgroup

La tabla fija una regla de lectura: una constante queda definida por su
objeto, su operador y su mapa de extracci\'on, no por la aparici\'on de un
prefijo decimal en un cat\'alogo. Euler--Mascheroni y Stieltjes pertenecen a
las partes finitas y los jets de Laurent; Ap\'ery, a una evaluaci\'on;
Glaisher--Kinkelin, a una derivada regularizada; Euler--Kronecker, al car\'acter
del campo; Meissel--Mertens, a una traza prima; Feigenbaum, a una familia
din\'amica finita; Catalan y Khinchin, a operadores vecinos distintos.

\subsection{Dominios, premisas y condiciones necesarias de la clasificación genealógica del número}

Las definiciones aritm\'eticas de entero, racional, irracional, algebraico y
trascendente son est\'andar. La carta ternaria balanceada, la separaci\'on entre
secci\'on y valor, el selector nativo de irreducibles, la ra\'iz Paley--Witt de
menos uno, las seis realizaciones de \(-1/12\) y los lectores espectrales se
deducen de los mapas citados en cada apartado. Esta secci\'on los reúne sin
introducir una constante ni un selector num\'erico adicional.

La fuerza de cada afirmaci\'on se determina por su mapa:
\begin{enumerate}
\item las biyecciones, identidades matriciales, coproductos, trazas y
      evaluaciones indicadas son exactas en sus dominios declarados;
\item la trascendencia de \(e\) y \(\pi\), la irracionalidad de
      \(\zeta(3)\) y el teorema erg\'odico de Khinchin son teoremas cl\'asicos
      incorporados como control o clasificaci\'on;
\item los aproximantes de Feigenbaum demuestran la rama y la profundidad
      calculadas; la universalidad requiere el operador de renormalizaci\'on y
      su teorema de convergencia;
\item la igualdad de dos escalares procedentes de dominios diferentes
      autoriza una comparaci\'on y exige un entrelazador para identificar los
      objetos completos;
\item una truncaci\'on decimal demuestra el prefijo calculado, mientras la
      clase aritm\'etica del valor infinito exige, según el caso, una ecuación
      polin\'omica o un teorema de irracionalidad o de trascendencia.
\end{enumerate}

Con estas reglas, la ontolog\'ia HMT del n\'umero adquiere una formulaci\'on
precisa: el valor arquimediano es la coordenada publicada; la secci\'on
compatible conserva la cadena que la genera; el operador determina qu\'e
propiedad se lee; y sus hipótesis determinan exactamente qu\'e conclusión se
deduce.


\section[Cierre genealógico del concepto de número]{Cierre genealógico del concepto de número: compatibilidad, memoria y evaluación arquimediana}
\label{p3-83-c44-s2:chap:cierre-genealogico-numero}
\index{número!cierre genealógico}
\index{cifra!coordenada visible}
\index{valor!evaluación arquimediana}
\index{conservación evolutiva de la unidad}

La teoría de números desarrollada en HMT culmina en una distinción de tipos
que gobierna la transición posterior hacia Mecánica Dimensional. Una
\emph{cifra} es una coordenada visible de un estado finito; un \emph{valor}
es la imagen arquimediana publicada por un lector; un \emph{número HMT} es
la genealogía compatible completa que conserva, junto con sus prefijos,
orientación, cociente, acarreo, frontera y memoria. La fibra de un valor
reúne las genealogías distintas que publican ese mismo escalar.

Esta construcción amplía la clasificación aritmética ordinaria. Entero,
racional, irracional, algebraico y trascendente siguen designando
propiedades del valor. HMT incorpora un segundo eje: la forma en que una
sección se prolonga, retorna, conserva memoria y alcanza una evaluación
estable. El cierre de las Partes~I--III fija ambos ejes y el mapa que los relaciona.

\subsection{Cifra, sección, valor y simetría}

Sea
\[
 \mathcal S_0^{\rm surv}
 \xleftarrow{\rho_{1,0}}
 \mathcal S_1^{\rm surv}
 \xleftarrow{\rho_{2,1}}
 \cdots
\]
el sistema inverso de estados supervivientes, donde cada restricción
conserva los datos tipados del TPK que pertenecen al nivel anterior, y sea
\[
 \mathcal S_\infty^{\rm surv}
 :=\varprojlim_n\mathcal S_n^{\rm surv}.
\]

\begin{definition}[Cifra]
\label{p3-83-c44-s2:def:cifra-visible-genealogica}
Fijado un lector posicional \(L\), la cifra o bloque visible a profundidad
\(n\) es
\[
 d_{L,n}(x_n):=\pi^{\rm vis}_{L,n}(x_n).
\]
Es la coordenada publicada por la presentación finita \(x_n\); la fibra de
\(\pi^{\rm vis}_{L,n}\) conserva los estados que comparten esa coordenada.
\end{definition}

\begin{definition}[Número HMT y presentación]
\label{p3-83-c44-s2:def:numero-hmt-cierre-genealogico}
Un número HMT es una sección compatible
\[
 \boldsymbol x=(x_n)_{n\geq0}\in\mathcal S_\infty^{\rm surv},
 \qquad
 \rho_{n+1,n}(x_{n+1})=x_n.
\]
El par \((\boldsymbol x,L)\) es una presentación del número. La sección
conserva la historia completa; el lector declara qué coordenadas de esa
historia se publican.
\end{definition}

\begin{definition}[Valor arquimediano]
\label{p3-83-c44-s2:def:valor-arquimediano-cierre-genealogico}
Supóngase que \(L\) asigna a cada \(x_n\) un cilindro compacto
\(I_{L,n}(\boldsymbol x)\subset\mathbb R\), con
\[
 I_{L,n+1}(\boldsymbol x)\subseteq I_{L,n}(\boldsymbol x),
 \qquad
 \operatorname{diam}I_{L,n}(\boldsymbol x)\longrightarrow0.
\]
El valor publicado por la presentación es
\[
 \operatorname{Val}_L(\boldsymbol x)
 :=\text{el único elemento de }
 \bigcap_{n\geq0}I_{L,n}(\boldsymbol x).
\]
\end{definition}

\begin{proposition}[Existencia y unicidad de la evaluación]
\label{p3-83-c44-s2:prop:evaluacion-seccion-genealogica}
Las hipótesis anteriores determinan una aplicación
\[
 \operatorname{Val}_L:\mathcal S_\infty^{\rm surv}\longrightarrow\mathbb R.
\]
Cada fibra \(\operatorname{Val}_L^{-1}(r)\) conserva las genealogías que la
evaluación arquimediana representa mediante el mismo valor \(r\).
\end{proposition}

\begin{proof}
La propiedad de intersección finita de la familia encajada de compactos
garantiza una intersección no vacía. Dos puntos distintos tendrían una
distancia positiva acotada por todos los diámetros, en contradicción con la
convergencia de éstos a cero.
\end{proof}

\begin{definition}[Simetría genealógica]
\label{p3-83-c44-s2:def:simetria-genealogica-hmt}
Una simetría del sistema de números HMT es una familia de biyecciones
\(g=(g_n)_{n\geq0}\) que preserva supervivencia y tipos, y satisface
\[
 \rho_{n+1,n}\circ g_{n+1}
 =g_n\circ\rho_{n+1,n}
 \qquad(n\geq0).
\]
La acción
\[
 g_\infty(\boldsymbol x):=(g_n x_n)_{n\geq0}
\]
es un automorfismo del límite inverso. Un invariante genealógico es una
función constante sobre sus órbitas.
\end{definition}

La simetría queda así caracterizada como transformación compatible a toda
profundidad. Esta noción hace comparables la publicación decimal, la
incidencia excepcional y la realización espectral posterior, conservando el
tipo propio de cada salida.

\subsection{Conservación evolutiva de la unidad}

Para \(n\geq1\), sean
\begin{equation}
 \rho_9(n)=1+((n-1)\bmod 9),
 \qquad
 \kappa_9(n)=\frac{n-\rho_9(n)}9.
 \label{p3-83-c44-s2:eq:residuo-cociente-conservacion-unidad}
\end{equation}
Entonces
\[
 n=9\kappa_9(n)+\rho_9(n),
 \qquad
 \rho_9(n)\in\{1,\ldots,9\}.
\]
La coordenada \(\rho_9\) publica la fase nonádica y representa la clase nula
mediante \(9\); \(\kappa_9\) conserva el número de vueltas. La raíz digital
publica la fase. La reducción módulo nueve acompañada de su cociente publica
fase y memoria.

La realización helicoidal asociada es
\begin{equation}
 \mathcal H_9(n):=
 \left(
 \cos\frac{2\pi(\rho_9(n)-1)}9,
 \sin\frac{2\pi(\rho_9(n)-1)}9,
 \kappa_9(n)+\frac{\rho_9(n)-1}{9}
 \right).
 \label{p3-83-c44-s2:eq:helice-nonadica-conservacion-unidad}
\end{equation}

\begin{proposition}[Retorno de fase y avance de memoria]
\label{p3-83-c44-s2:prop:retorno-fase-avance-memoria}
La realización \eqref{p3-83-c44-s2:eq:helice-nonadica-conservacion-unidad} satisface
\[
 \mathcal H_9(n+9)=\mathcal H_9(n)+(0,0,1).
\]
El retorno visible \(9\to1\) cierra la fase y eleva en una unidad la memoria
del estado.
\end{proposition}

\begin{proof}
Las identidades
\(\rho_9(n+9)=\rho_9(n)\) y
\(\kappa_9(n+9)=\kappa_9(n)+1\) fijan las coordenadas circulares e
incrementan exactamente la altura.
\end{proof}

Esta identidad formula la conservación evolutiva: la unidad retorna como
fase y el estado nuevo transporta el cierre anterior como memoria. La
completación cúbica realiza la misma ley en dimensión tres:
\begin{equation}
 10^3=(9+1)^3
 =9^3+3\cdot9^2+3\cdot9+1
 =729+243+27+1
 =729+271.
 \label{p3-83-c44-s2:eq:emrh-conservacion-evolutiva-unidad}
\end{equation}
El interior posee \(729\) estados; \(243+27\) constituye la corona
perforada; el ápice \(+1\) completa la unidad de base mil. La identidad
\(729+270=999\) describe el estado inmediatamente anterior al cierre, y
\(729+271=1000\) describe la completación. La ley
\[
 10^d=(9+1)^d
 =\sum_{r=0}^{d}\binom dr9^{d-r}
\]
transporta la misma unidad entre dimensiones y conserva la genealogía de sus
estratos.

\subsection{Sintaxis visible y semántica de prolongación}

Una firma local describe la forma actual de una rama. Su semántica de
supervivencia es el sistema de prolongaciones que puede sostener. Para
\(x_n\in\mathcal S_n^{\rm surv}\), defínase
\[
 \operatorname{Ext}_m(x_n)
 :=\{x_m\in\mathcal S_m^{\rm surv}:\rho_{m,n}(x_m)=x_n\},
 \qquad m\geq n.
\]
Dos estados que satisfacen
\[
 d_{L,n}(x_n)=d_{L,n}(y_n)
\]
pueden poseer árboles de extensión diferentes. En el testigo finito
\(60\to81\), dos ramas con la misma firma visible poseen, respectivamente,
cero y doce prolongaciones. El testigo establece que la continuidad futura
añade información a la salida local. El número HMT incorpora esa continuidad
como parte del objeto.

La coinducción se expresa mediante la compatibilidad de todos los prefijos.
Los bloques de mil o diez mil cifras son publicaciones truncadas de una regla
uniformemente prolongable. La generación numérica alcanza cualquier
profundidad solicitada; el levantamiento al estado enriquecido conserva,
además, las hipótesis de existencia, unicidad y compatibilidad de sus fibras.

\subsection{Clasificación aritmética y clasificación genealógica}

En una base \(b\geq2\), la presentación canónica de un real separa las clases
visibles siguientes:
\begin{enumerate}
\item un entero pertenece a \(\mathbb Z\) y admite una palabra ternaria
      balanceada finita;
\item un racional posee una expansión finita o finalmente periódica;
\item un irracional posee una expansión no finalmente periódica;
\item un algebraico es raíz de un polinomio no nulo de \(\mathbb Q[X]\);
\item un trascendente pertenece al complemento de la clausura algebraica de
      \(\mathbb Q\) en \(\mathbb C\).
\end{enumerate}
La periodicidad distingue racionales e irracionales; la ecuación polinómica
distingue algebraicos y trascendentes. Son ejes diferentes. Por ejemplo,
\(\varphi\) es irracional algebraico porque
\(\varphi^2-\varphi-1=0\); la trascendencia de \(e\) y \(\pi\) procede de
teoremas clásicos; una genealogía HMT distingue, adicionalmente, las
funciones de propagación, cierre y autoescala que realizan esos valores.

La misma separación rige para \(e+\pi\). La suma de las dos presentaciones
define un objeto compuesto y publica
\[
 \operatorname{Val}(\boldsymbol e)+
 \operatorname{Val}(\boldsymbol\pi)=e+\pi.
\]
Su clase aritmética constituye una pregunta propia. La construcción HMT
determina el mapa que forma el compuesto y conserva las dos genealogías que
intervienen.

\subsection{La familia operatorial de Euler como gramática de cierre}

La ecuación clásica de Euler pertenece a una familia cuadrática única. Sea
\(J_\sigma\) un generador que satisface
\[
 J_\sigma^2=\sigma I,
 \qquad \sigma\in\{-1,0,+1\}.
\]
La separación de las potencias pares e impares proporciona
\begin{equation}
 \exp(tJ_\sigma)=C_\sigma(t)I+S_\sigma(t)J_\sigma,
 \label{p3-83-c44-s2:eq:euler-trigeometrico-cierre-numero}
\end{equation}
con las realizaciones exactas
\[
\begin{array}{ccl}
J_{-1}^2=-I&:&\exp(tJ_{-1})=\cos t\,I+\sin t\,J_{-1},\\[1mm]
J_{0}^2=0&:&\exp(tJ_{0})=I+tJ_{0},\\[1mm]
J_{+1}^2=I&:&\exp(tJ_{+1})=\cosh t\,I+\sinh t\,J_{+1}.
\end{array}
\]
Elíptico, parabólico e hiperbólico son tres regímenes de una misma
exponencial operatorial. La clasificación trítica determina la geometría de
la evolución.

La rama hiperbólica contiene la autoescala de Fibonacci. Para
\[
 F=\begin{pmatrix}0&1\\1&1\end{pmatrix},
 \qquad
 A=F^2=\begin{pmatrix}1&1\\1&2\end{pmatrix},
\]
se tiene \(\det A=1\), \(\operatorname{tr}A=3\) y
\[
 \eta=\operatorname{arcosh}\frac32=2\log\varphi.
\]
Por tanto existe \(J_h^2=I\) tal que
\begin{equation}
 A=\exp(\eta J_h),
 \qquad
 A^n=\cosh(n\eta)I+\sinh(n\eta)J_h.
 \label{p3-83-c44-s2:eq:fibonacci-euler-hiperbolico-cierre-numero}
\end{equation}
La relación \(\varphi^2=\varphi+1\) expresa dos lecturas del mismo proceso:
recurrencia aditiva local y escala multiplicativa espectral. Esta
compatibilidad ofrece el modelo elemental de la doble evaluación
suma--producto que APP formaliza sobre su soporte común.

\subsection{Unificación tipada de las geometrías circular, compleja e hiperbólica}
\label{p3-83-c44-s2:sec:unificacion-compleja-circular-hiperbolica-numero}

La motivación geométrica inicial de HMT puede formularse exactamente como una
arquitectura común de soporte, frontera y generadores. El plano complejo
\(\mathbb C\) contiene el círculo unitario
\[
 S^1=\{e^{\mathrm i\theta}:\theta\in\mathbb R\},
\]
que constituye la realización elíptica de
\eqref{p3-83-c44-s2:eq:euler-trigeometrico-cierre-numero}. El semiplano superior
\(\mathbb H=\{w\in\mathbb C:\operatorname{Im}w>0\}\) se transporta al disco
hiperbólico mediante la transformación de Cayley
\[
 \mathcal C(w)=\frac{w-\mathrm i}{w+\mathrm i},
 \qquad
 \mathcal C:\mathbb H\xrightarrow{\sim}\mathbb D,
 \qquad
 \mathbb D=\{z\in\mathbb C:|z|<1\}.
\]
La frontera ideal de \(\mathbb D\) es el mismo círculo \(S^1\), mientras su
interior porta la métrica de Poincaré
\[
 \mathrm ds_{\mathbb D}^2
 =\frac{4\,|\mathrm dz|^2}{(1-|z|^2)^2}.
\]
La fusión consiste, por tanto, en un diagrama tipado: la parametrización
trigonométrica publica la frontera circular; la estructura compleja
proporciona la carta; la métrica de Poincaré realiza el interior hiperbólico.
Los tres regímenes \(J^2=-I,0,+I\) incorporan además el umbral parabólico y
sitúan esta concurrencia geométrica dentro de una sola familia operatorial.

\begin{proposition}[Compatibilidad euclídea e hiperbólica del cuarto de giro]
\label{p3-83-c44-s2:prop:cuarto-giro-poincare-rev3}
Sobre el disco unitario
\[
 \mathbb D=\{z\in\mathbb C:|z|<1\},
\]
la aplicación
\[
 \mathcal R_{\pi/2}:\mathbb D\longrightarrow\mathbb D,
 \qquad
 \mathcal R_{\pi/2}(z)=\mathrm i z,
\]
es simultáneamente una rotación euclídea y una isometría de la métrica de
Poincaré
\[
 \mathrm ds_{\mathbb D}^{2}
 =\frac{4\,|\mathrm dz|^2}{(1-|z|^2)^2}.
\]
Su extensión al borde \(S^1=\partial\mathbb D\) es el cuarto de giro
orientado. En la realización real del
\cref{thm:dos-realizaciones-fuente-cuadratica-rev3}, esta aplicación
corresponde a
\(\operatorname{Exp}((\pi/2)J_{\mathrm e})=J_{\mathrm e}\).
\end{proposition}

\begin{proof}
La multiplicación por \(\mathrm i\) conserva \(|z|\) y
\(|\mathrm dz|\). Conserva, por tanto, tanto la métrica euclídea como el
numerador y el denominador de \(\mathrm ds_{\mathbb D}^{2}\). La identidad
exponencial resulta de \(J_{\mathrm e}^2=-I\).
\end{proof}

El plano complejo, el círculo trigonométrico y el disco hiperbólico no se
identifican como espacios métricos: comparten un portador coordinado y el
automorfismo de cuarto de giro después de declarar cada realización.

\subsection{Interfaz de frontera entre la rotación cuadrática y la torsión mínima}
\label{p3-83-c44-s2:sec:interfaz-frontera-rotacion-torsion-minima-rev3}

La función de esta interfaz es transportar una rotación interior ya
construida hacia un dato de borde sin confundir el endomorfismo cuadrático
con un defecto escalar. Fijemos un punto base
\(x_0\in\partial X\), un fibrado real
\(E_{\partial}\to\partial X\), una conexión \(\nabla_{\partial}\), un
lector de frontera \(q_{\partial}\) y una co-tétrada discreta
\(e_{\partial}\in\Omega^1(\partial X;E_{\partial})\). Escribimos
\[
 \operatorname{Tors}(\partial X;E_{\partial})
 :=\Omega^2(\partial X;E_{\partial}),
 \qquad
 \operatorname{Hol}_{x_0}(\nabla_{\partial})
 \subseteq\operatorname{Aut}(E_{\partial,x_0}).
\]
Una realización de borde de la interfaz deberá construir los mapas
\[
 \operatorname{Res}_{\partial}:
 \langle J_{\mathrm e}\rangle
 \longrightarrow\operatorname{Hol}_{x_0}(\nabla_{\partial}),
 \qquad
 \operatorname{Def}_{\partial}:
 \operatorname{Hol}_{x_0}(\nabla_{\partial})
 \longrightarrow\operatorname{Tors}(\partial X;E_{\partial}).
\]
La torsión determinada por la conexión y la co-tétrada tiene el tipo
\begin{equation}
 T_{\partial}
 :=d_{\nabla_{\partial}}e_{\partial}
 \in\Omega^2(\partial X;E_{\partial}).
 \label{p3-83-c44-s2:eq:torsion-borde-interfaz-euler-rev3}
\end{equation}
La realización física de la torsión mínima requiere además un funcional
\begin{equation}
 \Lambda_{\partial}:
 \operatorname{Tors}(\partial X;E_{\partial})\longrightarrow\mathbb R
 \label{p3-83-c44-s2:eq:funcional-borde-delta4-interfaz-rev3}
\end{equation}
tal que los tres mapas satisfagan las condiciones de entrelazamiento
\begin{equation}
 \operatorname{Def}_{\partial}
 \bigl(\operatorname{Res}_{\partial}(J_{\mathrm e})\bigr)=T_{\partial},
 \qquad
 \Lambda_{\partial}(T_{\partial})=\Delta_4.
 \label{p3-83-c44-s2:eq:condiciones-borde-delta4-interfaz-rev3}
\end{equation}
Estas ecuaciones especifican el comparador que debe construirse; no afirman
que los tres mapas existan. El escalar \(\Delta_4\) está definido por
\cref{eq:delta4-electron-rev6}.

\begin{remark}[Condición exacta del comparador de frontera]
La raíz operatorial \(J_{\mathrm e}^2=-I\), el cuarto de giro
\(\operatorname{Exp}((\pi/2)J_{\mathrm e})=J_{\mathrm e}\) y el escalar
\(\Delta_4\) se deducen en sus dominios respectivos. La cadena de borde
anterior fija el tipo de los morfismos que deben enlazarlos. Mientras no se
construyan
\(\operatorname{Res}_{\partial}\), \(\operatorname{Def}_{\partial}\) y
\(\Lambda_{\partial}\), las condiciones
\eqref{p3-83-c44-s2:eq:condiciones-borde-delta4-interfaz-rev3} definen el comparador
requerido pero no implican una realización física; la proximidad angular de
sus extremos no sustituye ese operador.
\end{remark}

APP, TRIT y TPK suministran el soporte discreto de esa arquitectura. Sobre
\(X_9=(\mathbb Z/9\mathbb Z)^2\), las evaluaciones suma y producto inducen
proyectores \(P_{\Sigma,d}\) y \(P_{\Pi,d}\) que satisfacen
\[
 \|[P_{\Sigma,2},P_{\Pi,2}]\|=\frac{\sqrt2}{3},
 \qquad
 \|[P_{\Sigma,3},P_{\Pi,3}]\|=\frac{\sqrt3}{4}.
\]
Esta incompatibilidad calculada es el precursor finito de Heisenberg: dos
familias de observables comparten soporte y poseen descomposiciones
espectrales incompatibles. La forma alternante del toro,
\[
 \sigma((a,b),(c,d))=ad-bc\pmod9,
\]
se integra en la representación nonádica de Weyl--Heisenberg
\[
 X|n\rangle=|n+1\rangle,\qquad
 Z|n\rangle=\omega^n|n\rangle,\qquad
 ZX=\omega XZ,
 \qquad \omega=e^{2\pi\mathrm i/9}.
\]
El conmutador publica la fase orientada y el grupo central resultante posee
\(9^3=729\) elementos.

El bloque central de APP proporciona simultáneamente la realización
lorentziana local:
\[
 B_c=
 \begin{pmatrix}7&2\\2&7\end{pmatrix}
 =9P_++5P_-,
 \qquad
 M_c:=\frac{B_c}{\sqrt{45}}
 =\exp\!\left(\frac12\log\frac95\,R\right).
\]
Con \(\eta_{1,1}=\operatorname{diag}(1,-1)\),
\[
 M_c^{\mathsf T}\eta_{1,1}M_c=\eta_{1,1},
 \qquad
 \det M_c=1,
 \qquad
 M_c\in SO^+(1,1).
\]
Así, la geometría de Lorentz y el plano de Minkowski \(1+1\) emergen como
realización espectral exacta del bloque central. Toda elevación dimensional
posterior debe entrelazar este bloque local con su nuevo codominio. La secuencia
\[
 \boxed{
 \text{APP--TRIT--TPK}
 \longrightarrow
 \begin{cases}
 \text{familia elíptica--parabólica--hiperbólica},\\
 \text{bloque Lorentz--Minkowski }1+1,\\
 \text{representación finita de Weyl--Heisenberg}
 \end{cases}}
\]
expone los desarrollos geométricos y operatoriales que preceden al cierre
genealógico del número.

\subsection{Operadores de generación de constantes: dominios, semillas y evaluaciones}

La función de una constante se fija mediante cuatro datos: definición
matemática, dominio, operador o lector y conclusión deducida. La
\cref{p3-83-c44-s2:tab:semillas-cierre-genealogico-numero} reúne los objetos que deben
permanecer visibles al cerrar la teoría de números.

\begingroup
\small
\setlength{\LTpre}{0.5\baselineskip}
\setlength{\LTpost}{0.5\baselineskip}
\begin{longtable}{@{}>{\raggedright\arraybackslash}p{.145\textwidth}
                    >{\raggedright\arraybackslash}p{.235\textwidth}
                    >{\raggedright\arraybackslash}p{.34\textwidth}
                    >{\raggedright\arraybackslash}p{.20\textwidth}@{}}
\caption{Constantes y semillas en el cierre genealógico del número.}
\label{p3-83-c44-s2:tab:semillas-cierre-genealogico-numero}\\
\toprule
\textbf{Objeto} & \textbf{Definición o identidad} &
\textbf{Dominio y mapa HMT} & \textbf{Conclusión y condición}\\
\midrule
\endfirsthead
\toprule
\textbf{Objeto} & \textbf{Definición o identidad} &
\textbf{Dominio y mapa HMT} & \textbf{Conclusión y condición}\\
\midrule
\endhead
\(\pi,e,\varphi\) &
Secciones publicadas por cilindros numéricos compatibles. &
La filtración
\(w_6\to w_{30}\to R_{36}\to G_9\) produce el estado terminal; los lectores
separan cierre, propagación y autoescala. &
Para cada profundidad finita prescrita, el mapa produce cilindros unitarios
compatibles; la sección infinita resulta del teorema coinductivo.\\
\(\alpha_{\rm HMT}\) &
Cierre dodecafásico con acarreo sobre el estado terminal. &
La doble lectura conduce a
\((P,E,\Phi,K,c_{13}=0)\) y a la salida racional de \(36\) cifras; la vía
\(E_{21/22}\) proporciona un cálculo independiente. &
La recurrencia es exacta para el estado terminal declarado y no recibe
\(\alpha\) ni CODATA; la deducción mínima APP\(\to K\) requiere construir su
lector bidireccional completo.\\
\(i\) interno &
\(J^2=-I\) sobre el módulo Paley--Witt. &
Realiza una estructura compleja finita y separa orientaciones conjugadas. &
Sobre el cuerpo base declarado se cumple exactamente \(J^2=-I\).\\
\(-1/12\) &
\(\zeta(-1)=-1/12\) y realizaciones operatoriales normalizadas. &
Incidencia \(90\to120\), defecto dodecafásico y funcionales espectrales,
cada uno con su dominio y normalización. &
Igualdad exacta por cada mapa declarado.\\
\(e+\pi\) &
Suma de las presentaciones de \(e\) y \(\pi\). &
Objeto compuesto del álgebra de presentaciones; conserva ambas genealogías. &
Construcción exacta; clase aritmética separada.\\
Euler--Mascheroni \(\gamma\) y Stieltjes \(\gamma_n\) &
\(\gamma=\operatorname*{FP}_{s=1}\zeta(s)\); coeficientes de Laurent. &
Descomposición residual módulo tres y jets de Laurent. &
Identidades exactas del funcional declarado.\\
Apéry \(\zeta(3)\) &
Evaluación de la zeta en \(s=3\). &
Lector \(\mathcal Z_3(3)\) con intervalos racionales de prefijo. &
Evaluación exacta; irracionalidad clásica.\\
Glaisher--Kinkelin \(A_{\rm GK}\) &
\(A_{\rm GK}=\exp(1/12-\zeta'(-1))\). &
Derivada regularizada del lector espectral. &
Identidad exacta con normalización explícita.\\
Euler--Kronecker \(\gamma_{F_{-3}}\) &
Parte finita logarítmica de la zeta de Dedekind de
\(F_{-3}=\mathbb Q(\sqrt{-3})\). &
El carácter \(\chi_{-3}\) selecciona el campo y organiza
\(L'(1,\chi_{-3})/L(1,\chi_{-3})\). &
Identidad exacta del carácter declarado.\\
Meissel--Mertens \(B_1\) &
\(B_1=\lim_{x\to\infty}(\sum_{p\le x}p^{-1}-\log\log x)\). &
Parte finita de la traza sobre modos primos. &
Representación operatorial exacta con modos fijados.\\
Catalan \(G\) y Khinchin \(K_0\) &
\(G=L(2,\chi_{-4})\); media ergódica del sistema de Gauss. &
Lectores de carácter módulo cuatro y de fracción continua. &
Teoremas clásicos incorporados con dinámicas propias.\\
Feigenbaum \((\delta,\alpha_F)\) &
Escalas del operador de renormalización en duplicación de periodo. &
La familia dodecafásica produce aproximantes de nivel ocho. &
El cálculo determina exactamente el nivel ocho; la universalidad requiere
demostrar la convergencia al punto fijo de renormalización.\\
Rogers--Ramanujan \(R(q)\) &
Fracción continua de nivel cinco e identidad de Klein--Ramanujan para
\(j\). &
Organiza nivel \(5\), \(A_5\), el límite
\(R(q)\to\varphi^{-1}\) y la replicabilidad Hecke--Faber. &
Identidades modulares exactas; la comparación de coeficientes sólo cubre la
profundidad finita calculada.\\
Reloj primo \(p\) &
\(\tau_p=\operatorname{ord}_p(10)\),
\(M_p=(10^{\tau_p}-1)/p\equiv0\pmod9\) para \(p\nmid30\). &
Los primos son periodos irreducibles; los compuestos sincronizan relojes
mediante el mínimo común múltiplo. &
Teorema aritmético exacto en el dominio declarado.\\
\bottomrule
\end{longtable}
\endgroup

Parte finita, valor especial, derivada regularizada, carácter, promedio
ergódico, escala de renormalización y periodo primo son mapas distintos. Su
convivencia dentro de HMT procede de una familia tipada de operadores y
mapas de extracción.

\subsection{Doble proyección de una misma genealogía}

Sea \(\mathfrak X:=X_\infty^{\rm cal}\) el espacio compacto de historias
enriquecidas con calibre inicial transportado. Fijado un calibre declarado
\(c\in\mathscr C_{\rm exc}\), sobre una misma sección actúan dos
evaluaciones:
\[
 \mathfrak R:\mathfrak X\longrightarrow\mathbb R,
 \qquad
 \mathfrak E_c:\mathfrak X\longrightarrow\mathbf{Exc},
 \qquad
 \mathfrak E_c(x):=\mathfrak E(x,c),
\]
donde \(\mathfrak R\) contrae cilindros y \(\mathfrak E_c\) conserva soportes,
incidencias, polaridad y transporte de frontera.

\begin{theorem}[Unidad causal de la doble proyección]
\label{p3-83-c44-s2:thm:unidad-causal-doble-proyeccion-numero}
Para toda historia \(\boldsymbol x\in\mathfrak X\), el par
\[
 \mathfrak D_c(\boldsymbol x)
 :=(\mathfrak R(\boldsymbol x),\mathfrak E_c(\boldsymbol x))
\]
es una publicación conjunta del mismo antecedente. La primera coordenada
determina el valor arquimediano; la segunda determina la clase de incidencia
excepcional. Cuando los lectores conjuntos separan puntos,
\(\mathfrak D_c\) es inyectiva.
\end{theorem}

\begin{proof}
Ambas componentes están definidas sobre el mismo dominio y conservan los
mapas de restricción. Si
\(\mathfrak D_c(\boldsymbol x)=\mathfrak D_c(\boldsymbol y)\), la separación
conjunta de historias implica \(\boldsymbol x=\boldsymbol y\).
\end{proof}

La rama arquimediana produce los cilindros y los valores de
\(\pi,e,\varphi\). La rama de incidencia prolonga las mismas palabras y el
mismo sello hacia Paley, Witt, Golay, el retículo de Leech,
\(V^\natural\), el grupo Monstruo y Moonshine. La relación causal es la
preimagen común: ambas evaluaciones conservan invariantes diferentes de la
misma historia. El real es la sombra contraída; la simetría excepcional es la
memoria de incidencia publicada por el mismo estado.

\subsection{Clausura genealógica del concepto de número}

El concepto de número queda fijado por la cadena
\[
 \boxed{
 \text{estado finito}
 \longrightarrow\text{prolongación compatible}
 \longrightarrow\text{sección genealógica}
 \longrightarrow
 \begin{cases}
   \text{cifra visible},\\
   \text{valor arquimediano},\\
   \text{incidencia excepcional}.
 \end{cases}}
\]
La cifra es una coordenada; el valor es una proyección; el número es la
sección completa; la simetría es el automorfismo compatible que conserva su
memoria. La conservación evolutiva transporta esa identidad a través del
retorno nonádico, del acarreo helicoidal y de la completación
\(729+271=1000\). Con estas distinciones cerradas, Mecánica Dimensional
realiza la misma genealogía en un nuevo codominio.

\paragraph{Dependencias matemáticas.}
El límite inverso proporciona la sección; el lector arquimediano produce el
valor; la completación binomial y la extensión nonádica transportan unidad y
memoria; los lectores complementarios clasifican las salidas; y la doble
proyección aplica al mismo antecedente los mapas arquimediano y excepcional.
La cadena Witt--Golay--Leech--Moonshine requiere, en cada paso, el
entrelazador de incidencia correspondiente.
